



































































































































































































































































































































































































































































































































































\section{Introduction}
For a (labeled) graph $G=(V,E)$, and a bijection $\phi:V\to V'$,
we denote by $\phi(G)$ the graph $G'=(V',E')$
such that $E'=\{\{\phi(u),\phi(v)\}:\{u,v\}\!\in\!E\}$,
and say that $G'$ is {\sf isomorphic} to~$G$.
The set of {\sf automorphisms} of the graph $G=(V,E)$,
denoted ${\tt aut}(G)$,
is the set of permutations that preserve the graph~$G$;
that is, $\pi\in{\tt aut}(G)$ if and only if $\pi(G)=G$.
We say that a graph is {\sf asymmetric} (equiv., {\sf self-ordered})
if its set of automorphisms is a singleton,
which consists of the trivial automorphism
(i.e., the identity permutation).
We actually prefer the term {\em self-ordered}, because we take
the perspective that is offered by the following equivalent definition.

%
\BD[Self-ordered (a.k.a asymmetric) graphs]
\label{asymmetric:def}
%
The graph $G=([n],E)$ is {\sf self-ordered}
if for every graph $G'=(V',E')$ that is isomorphic to~$G$
there exists a unique bijection $\phi:V'\to[n]$ such that $\phi(G')=G$.
\ED
In other words,
given an isomorphic copy $G'=(V',E')$ of a fixed graph $G=([n],E)$,
there is a unique bijection $\phi:V'\to[n]$ that orders the vertices of $G'$
such that the resulting graph (i.e., $\phi(G')$) is identical to~$G$.
Indeed, if $G'=G$, then this unique bijection is the identity permutation.
\footnote{Naturally, we are interested in efficient algorithms
that find this unique ordering, whenever it exists;
such algorithms are known
when the degree of the graph is bounded~\cite{L}.}

In this work, we consider a feature,
which we call {\em robust self-ordering},
that is a quantitative version self-ordering.
%
Loosely speaking, a graph $G=([n],E)$ is robustly self-ordered if,
%
%
for every permutation $\pi:[n]\to[n]$,
the size of the symmetric difference between~$G$ and $\pi(G)$
is proportional to the number of non-fixed-points under $\pi$;
that is, $|E\triangle\{\{\pi(u),\pi(v)\}\!:\!\{u,v\}\!\in\!E\}|$
is proportional to $|\{i\!\in\![n]\!:\!\pi(i)\!\neq\!i\}|$.
(In contrast, self-ordering only means that
the size of the symmetric difference is positive
if the number of non-fixed-points is positive.)

%
\BD[Robustly self-ordered graphs]
\label{robust-asymmetric:def}
A graph $G=(V,E)$ is said to be {\sf $\gamma$-robustly self-ordered}
if for every permutation $\pi:V\to{V}$
it holds that
\begin{equation}\label{robust-asymmetric:eqdef}
\big|E\triangle\left\{\{\pi(u),\pi(v)\}\!:\!\{u,v\}\!\in\!E\right\}\big|
  \;\geq\; \gamma\cdot|\{i\in[n]\!:\!\pi(i)\neq i\}|,
\end{equation}
where $\triangle$ denotes the symmetric difference operation.
An infinite family of graphs $\{G_n=([n],E_n)\}_{n\in\N}$
{\em(such that each $G_n$ has maximum degree $d$)}
is called {\sf robustly self-ordered}
if there exists a constant $\gamma>0$, called the {\sf robustness parameter},
such that for every $n$ the graph $G_n$ is $\gamma$-robustly self-ordered.
\ED
Note that
$|E_n\triangle\{\{\pi(u),\pi(v)\}\!:\!\{u,v\}\!\in\!E_n\}|
 \leq 2d\cdot|\{i\in[n]\!:\!\pi(i)\neq i\}|$
always holds (for families of maximum degree $d$).
%
%
%
%
%
The term ``robust'' is inspired by
the property testing literature (cf.~\cite{RS}),
where it indicates that some ``parametrized violation''
is reflected proportionally in some ``detection parameter''.

The second part of Definition~\ref{robust-asymmetric:def}
is tailored for bounded-degree graphs, which will be our focus
in Section~\ref{edge-colored:sec}--6.
Nevertheless,
in Sections~7--10
we consider graphs
of unbounded degree and unbounded robustness parameters.
In this case, for a function $\rho:\N\to\R$,
we say that an infinite family of graphs $\{G_n=([n],E_n)\}_{n\in\N}$
is {\sf $\rho$-robustly self-ordered} if for every $n$
the graph $G_n$ is $\rho(n)$-robustly self-ordered.
Naturally, in this case, the graphs
must have $\Omega(\rho(n)\cdot n)$ edges.
\footnote{Actually, all but at most one vertex
must have degree at least $\rho(n)/2$.}
%
In Sections~7--9
we consider the case of $\rho(n)=\Omega(n)$.
















\setcounter{theorem}{2}\BT[Constructing robustly self-ordered bounded-degree graphs]
\label{main:ithm}
For all sufficiently large $d\in\N$,
there exist an infinite family of $d$-regular
robustly self-ordered graphs $\{G_n\}_{n\in\N}$
and a polynomial-time algorithm that,
given $n\in\N$ and a vertex $v\in[n]$ in the $n$-vertex graph $G_n$,
finds all neighbors of $v$ {\rm(in $G_n$)}.
%
\ET
We stress that the algorithm runs in time that is polynomial
in the description of the vertex; that is, the algorithm runs
in time that is polylogarithmic in the size of the graph.
Theorem~\ref{main:ithm} holds both for graphs that consists
of connected components of logarithmic size
and for ``strongly connected'' graphs (i.e., expanders).






































































































































































































































































































































































































































































































































































































































































































































































































































































































































































































\section{The Edge-Colored Variant}
\label{edge-colored:sec}
Many of our arguments are easier to make in a model
of (bounded-degree) graphs in which edges are colored
(by a bounded number of colors),
and where one counts the number of mismatches between colored edges.
Namely, an edge that appears in one (edge-colored) graph
contributes to the count if it either does not appear in the
other (edge-colored) graph or appears in it under a different color.
Hence, we define a notion of robust self-ordering for edge-colored graphs.
We shall then transform robustly self-ordered edge-colored graphs
to robustly self-ordered ordinary (uncolored) graphs,
while preserving the degree, the asymptotic number of vertices,
and other features such as expansion and degree-regularity.
Specifically, the transformation consists of replacing the
colored edges by copies of different connected,
asymmetric (constant-sized) gadgets such that different colors
are reflected by different gadgets.

We start by providing the definition of the edge-colored model.
Actually, for greater flexibility, we will consider multi-graphs;
that is, graphs with possible parallel edges and self-loops.
Hence, we shall consider multi-graphs $G=(V,E)$
coupled with an edge-coloring function $\chi\!:\!E\!\to\!\N$,
where $E$ is a multi-set containing both pairs of vertices
and singletons (representing self-loops).
Actually, it will be more convenient to represent self-loops
as 2-element multi-sets containing two copies of the same vertex.

%
\BD[Robust self-ordering of edge-colored multi-graphs]
\label{colored-robust:def}
Let $G=(V,E)$ be a multi-graph with colored edges,
where $\chi\!:\!E\!\to\!\N$ denotes this coloring,
and let $E_i$ denote the multi-set of edges colored~$i$
{\rm(i.e., $E_i=\{e\!\in\!E\!:\!\chi(e)\!=\!i\}$)}.
We say that $(G,\chi)$ is {\sf $\gamma$-robustly self-ordered}
if for every permutation $\mu:V\to{V}$ it holds that
\begin{equation}\label{robust4colored:eqdef}
\sum_{i\in\N}
\Big|E_i\;\triangle
       \left\{\{\mu(u),\mu(v)\}\!:\!\{u,v\}\!\in\!E_i\right\}\Big|
  \;\geq\; \gamma\cdot|\{i\!\in\!V\!:\!\mu(i)\neq i\}|,
\end{equation}
where $A\triangle B$ denotes the symmetric difference
between the multi-sets $A$ and $B$; that is $A\triangle B$
contains $t$ occurrences of $e$ if the absolute difference between
the number of occurrences of $e$ in $A$ and $B$ equals~$t$.
\ED
(Definition~\ref{robust-asymmetric:def} is obtained
as a special case when the multi-graph is actually a graph
and all edges are assigned the same color.)

We stress that whenever we consider ``edge-colored graphs''
we actually refer to edge-colored multi-graphs
(i.e., we explicitly allow parallel edges and self-loops).
\footnote{We comment that a seemingly more appealing definition
can be used for edge-colored (simple) graphs.
Specifically, in that case (i.e., $E\subseteq\binom{V}{2}$),
we can extend $\chi\!:\!E\!\to\!\N$ to non-edges
by defining $\chi(\{u,v\})=0$ if $\{u,v\}\not\in E$,
and say that $(G,\chi)$ is {\sf $\gamma$-robustly self-ordered}
if for every permutation $\mu:V\to{V}$ it holds that\\[-1ex]
$$\left|\left\{\{u,v\}\in \binom{V}{2}:
       \chi(\{\mu(u),\mu(v)\})\!\neq\!\chi(\{u,v\})\right\}\right|
  \;\geq\; \gamma\cdot|\{i\!\in\!V:\mu(i)\neq i\}|.$$}
%
In contrast, whenever we consider (uncolored) graph,
we refer to simple graphs (with no parallel edges and no self-loops).

Our transformation of robustly self-ordered edge-colored multi-graphs
to robustly self-ordered ordinary graphs depends on the number
of colors used by the multi-graph.
In particular, $\gamma$-robustness of edge-colored multi-graph
that uses $c$ colors gets translated to $(\gamma/f(c))$-robustness
of the resulting graph, where $f:\N\to\N$ is an unbounded function.
Hence, we focus on coloring functions that use a constant number of colors,
denoted $c$. That is, fixing a constant $c\in\N$,
we shall consider multi-graphs $G=(V,E)$ coupled with an edge-coloring
function $\chi\!:\!E\!\to\![c]$.
%
%

\subsection{Transformation to standard (uncolored) version}
\label{edge-colored:trans:sec}
As a preliminary step for the transformation,
we add self-loops to all vertices
and make sure that parallel edges are assigned different colors.
The self-loops make it easy to distinguish
the original vertices from auxiliary vertices
that are parts of gadgets introduced in the main transformation.
Different colors assigned to parallel edges are essential
to the mere asymmetry of the resulting graph,
since we are going to replace edges of the same color
by copies of the same gadget.

\BCT[Preliminary step towards Construction~\ref{colored2standard:ct}]
\label{pre-colored2standard:ct}
For a fixed $d\geq3$,
given a multi-graph $G=(V,E)$ of maximum degree $d$
and an edge-coloring function $\chi\!:\!E\!\to\![c]$,
we define a multi-graph $G=(V,E')$ and
an edge-coloring function $\chi'\!:\!E'\!\to\![d\cdot c+1]$ as follows.
\BE
\item For every pair of vertices $u$ and~$v$ that are connected
by few parallel edges, denoted $e^{(1)}_{u,v},\ldots,e^{(d')}_{u,v}$,
we change, for each $i\in[d']$, the color of $e^{(i)}_{u,v}$
to $\chi'(e^{(i)}_{u,v})\gets(i-1)\cdot c+\chi(e^{(i)}_{u,v})$.
This includes also the case $u=v$.
\item We augment the multi-graph with self-loops colored $d\cdot c+1$;
that is, $E'$ is the multi-set $E\cup\{e_v:v\!\in\!V\}$,
where $e_v$ is a self-loop added to $v$, and $\chi'(e_v)=dc+1$.
\EE
{\rm(Other edges $e\!\in\!E$ maintain their color;
that is, for them $\chi'(e)=\chi(e)$ holds.)}
\ECT

(For simplicity, we re-color all parallel edges, save the first one,
rather than re-coloring only parallel edges that have the same color.)
Note that refining the coloring may only increase
the robustness parameter of an edge-colored multi-graph.
Clearly, $G'$ preserves many features of~$G$.
In particular, it preserves $\gamma$-robust self-ordering,
expansion, degree-regularity, and the number of vertices.

As stated above, our transformation of edge-colored multi-graphs
to ordinary graphs uses gadgets, which are constant-size graphs.
Specifically, when handling a multi-graph of maximum degree $d$
with edges that are colored by $c$ colors,
we shall use $c$ different {\em connected and asymmetric}\/ graphs.
Furthermore, in order to maintain $d$-regularity,
we shall use $d$-regular graphs as gadgets;
and in order to have better control on the number
of vertices in the resulting graph,
each of these gadgets will contain $k=k(d,c)$ vertices.
The existence of such ($d$-regular) asymmetric (and connected) graphs
is well-known, let alone that it is known that
a random $d$-regular $k$-vertex graph is asymmetric
(for any constant $d\geq3$)~\cite{B1,B2}.

We stress that the different gadgets are
each connected and asymmetric,
and it follows that they are not isomorphic to one another.
We designate in each gadget an edge $\{p,q\}$,
called the {\sf designated edge},
such that omitting this edge does not disconnect the gadget.
The endpoints of this edge
will be used to connect two vertices of the original multi-graph.
Specifically, we replace each edge $\{u,v\}$ (of the original multi-graph)
that is colored $i$ by a copy of the $i^\xth$ gadget,
while omitting its designated edge $\{p,q\}$,
and connecting $u$ to $p$ and $v$ to $q$.
The construction is spelled out below.

We say that a (non-simple) multi-graph $G=(V,E)$
coupled with an edge-coloring $\chi$ is {\sf eligible}
if each of its vertices contains a self-loop,
and parallel edges are assigned different colors.
Recall that eligibility comes almost for free
(by applying Construction~\ref{pre-colored2standard:ct}).
We shall apply the following construction
only to eligible edge-colored multi-graphs.

\BCT[The main transformation]
\label{colored2standard:ct}
For a fixed $d\geq3$ and $c$, let $k=k(d,c)$
and $G_1,\ldots,G_c$ be different asymmetric
and connected $d$-regular graphs over the vertex-set $[k]$.
%
Given a multi-graph $G=(V,E)$ of maximum degree $d$
and an edge-coloring function $\chi\!:\!E\!\to\![c]$,
we construct a graph $G'=(V',E')$ as follows.
\begin{quote}
%
Suppose that the multi-set $E$ has size $m$.
Then, for each $j\in[m]$, if the $j^\xth$ edge of $E$
connects vertices $u$ and~$v$, and is colored $i$,
then we replace it by a copy of $G_i$, 
%
while omitting its designated edge and connecting
one of its endpoints to $u$ and the other endpoint to $v$.

Specifically, assuming that $V=[n]$ and
recalling that $j$ is the index of the edge
{\rm(colored~$i$)} that connects $u$ and~$v$,
let $G_i^{u,v}$ be an isomorphic copy of $G_i$
that uses the vertex set $\{n+(j-1)\cdot k+i:i\!\in\![k]\}$.
Let $\{p,q\}$ be the designated edge in $G_i^{u,v}$,
and $\GG_i^{u,v}$ be the graph that results from $G_i^{u,v}$
by omitting $\{p,q\}$.
Then, we replace the edge $\{u,v\}$ by $\GG_i^{u,v}$,
and add the edges $\{u,p\}$ and $\{v,q\}$.
%
\end{quote}
Hence, $V'=[n+m\cdot k]$
and $E'$ consists of the edges of all $\GG_i^{u,v}$'s
as well as the edges connecting the endpoint of the corresponding
designated edges to the corresponding vertices~$u$ and~$v$.
\ECT
We stress that, although~$G$ may have parallel edges and self-loops,
the graph $G'$ has neither parallel edges nor self-loops.
Also note that $G'$ preserve various properties of~$G$
such as degree-regularity, number of connected components,
and expansion (up to a constant factor).

%
%
%
%
%
%
%
%
%
%
%

We shall show that if the edge-colored multi-graph $G=(V,E)$
is robustly self-ordered (in the edge-colored sense),
then the resulting graph $G'=(V',E')$ is robustly self-ordered
(in the ordinary sense). The proof of this fact relies on
a correspondence between the colored edges of~$G$ and the gadgets in $G'$.
For starters, suppose that the permutation $\mu':V'\to V'$
%
maps $V$ to $V$ (i.e., $\mu'(V)=V$),
and gadgets to the corresponding gadgets; that is,
if $\mu'$ maps the vertex-pair $(u,v)\in V^2$
to $(\mu'(u),\mu'(v))\in V^2$,
then $\mu'$ maps the vertices in the possible gadget that
connects~$u$ and~$v$ to the vertices in the gadget that
connects $\mu'(u)$ and $\mu'(v)$.
%
%
%
%
In such a case, letting $\mu$ be the restriction of $\mu'$ to $V$,
a difference of $D$ colored edges between~$G$ and $\mu(G)$
translates to a difference of
at least $D$ edges between~$G'$ and $\mu'(G')$,
due to the difference between the gadgets
that replace the corresponding (colored) edges of $G'$,
%
%
%
%
whereas the number of non-fixed-point vertices in $\mu'$ is $k$ times
larger than the number of non-fixed-point vertices in $\mu$.
%
Assuming that~$G$ is $\gamma$-robustly self-ordered,
it follows that $\mu$ has at most $D/\gamma$ non-fixed-points.
Hence, in this case we have
 \begin{equation*}
\frac{|G'\triangle\,\mu'(G')|}{|\{v\in V':\mu'(v)\!\neq\!v\}|}
%
  \geq \frac{D}{k\cdot|\{v\in V:\mu(v)\!\neq\!v\}|}
   \geq \frac{D}{k\cdot D/\gamma}
 \end{equation*}
which equals $\gamma/k$.
%
%
However, in general, $\mu'$ needs not satisfy the foregoing condition.
Nevertheless, if $\mu'$ splits some gadget
or maps some gadget in a manner that is inconsistent with
the vertices of $V$ connected by it,
then this gadget contributes at least one unit
to the difference between~$G'$ and $\mu'(G')$,
whereas the number of non-fixed-point vertices in this gadget
is at most~$k$.
Lastly, if $\mu'$ maps vertices of a gadget to other
vertices in the same gadget, then we get a contribution
of at least one unit due to the asymmetry of the gadget.
The foregoing argument is made rigorous in the proof of the following theorem.

\BT[From edge-colored robustness to standard robustness]
\label{colored2std:clm}
For constant $d\geq3$ and $c\in\N$,
suppose that the multi-graph $G=(V,E)$ coupled with $\chi\!:\!E\!\to\![c]$
is eligible and $\gamma$-robustly self-ordered.
%
%
%
%
%
%
%
%
%
%
%
Then, the graph $G'=(V',E')$ resulting from
Construction~\ref{colored2standard:ct}
is $(\gamma/3k)$-robustly self-ordered,
%
where $k=k(d,c)>d$ is the number of vertices in a gadget
{\rm(as determined above)} and $\gamma\leq1$.
%
%
\ET
%
%
%
%

\BPF
As a warm-up, let us verify that $G'$ is asymmetric.
We first observe that the vertices of~$G$ are uniquely identified
(in $G'$), since they are the only vertices that are incident at copies
of the gadget that replaces the self-loops.
%
%
\footnote{Indeed, a vertex of $G'$ is in $V$ if and only if
omitting it from $G'$ yields several connected components
such that (at least) one of them is a copy of the gadget
that replaces the self-loops (with the designated edge missing).}
%
Hence, any automorphism of $G'$ must map $V$ to $V$.
Consequently, for any $i$, such an automorphism $\mu'$
must map each copy of $G_i$ to a copy of $G_i$,
which means that when permuting $V$ according to $\mu'$
the edges of~$G$ as well as their colors are preserved.
%
By the ``colored asymmetry'' of~$G$,
this implies that $\mu'$ maps each $v\in V$ to itself,
and consequently each copy of $G_i$ must be mapped (by $\mu'$) to itself.
Finally, using the asymmetry of the $G_i$'s,
it follows that each vertex of each copy of $G_i$ is mapped to itself.
Hence, $\mu'$ must be the identity permutation.

We now turn to proving that $G'$ is actually robustly self-ordered.
%
%
%
%
%
%
%
Considering an arbitrary permutation $\mu':V'\to V'$,
we lower-bound the distance between~$G'$ and $\mu'(G')$
as a function of the number of non-fixed-points under $\mu'$
(i.e., of $v\in V'$ such that $\mu'(v')\neq v'$).
We do so by considering the contribution
of each non-fixed-point to the distance between~$G'$ and $\mu'(G')$.
We first recall the fact that the vertices of $V$
(resp., of gadgets) are uniquely identified in $\mu'(G')$
by virtue of the gadgets that replace self-loops
(see the foregoing warm-up).
%
%
%
%

\BDes
\item[{\em Case 1}:] {\em Vertices of some copy of $G_i$
that are not mapped by $\mu'$ to a single copy of $G_i$;
that is, vertices in some $G_i^{u,v}$ that are not
mapped by $\mu'$ to some $G_i^{u',v'}$.}

(This includes the case of vertices $w'$ and $w''$ of some $G_i^{u,v}$
such that $\mu'(w')$ is in $G_{i'}^{u',v'}$
and~$\mu'(w'')$ is in $G_{i''}^{u'',v''}$,
but $(i',u',v')\neq(i'',u'',v'')$.
It also includes the case of a copy of $G_i$ that is mapped by $\mu'$
to a copy of $G_j$ for $j\neq i$,
and the case that a vertex $w$ in some $G_i^{u,v}$
that is mapped by $\mu'$ to a vertex in $V$.)

%
%
%
%
The set of vertices $S_i^{u,v}$ of each such copy (i.e., $G_i^{u,v}$)
contribute at least one unit to the difference between~$G'$ and $\mu'(G')$,
since $\mu'(S_i^{u,v})$ induces a copy of $\GG_i$ in $\mu(G')$
but not in $G'$,
where here we also use the fact that the $\GG_i$'s are connected
(and not isomorphic (for the case of $i'=i''\neq i$)).
%
%
%
%
%
%
%
%
%
%
%
%
%
Note that the total contribution of all vertices of the current case
equals at least the number of gadgets in which they reside.
Hence, if the current case contains $n_1$ vertices,
then their contribution to the distance between~$G'$ and $\mu'(G')$
is at least $n_1/k$.

Ditto for vertices that do not belong to a single copy of $G_i$
and are mapped by $\mu'$ to a single copy of $G_i$.
This also includes $v\in V$ being mapped to some copy of some $G_i$,
but in this case we get a contribution of one unit
(rather than $1/k$ amortized units)
per each such vertex (i.e., $v\in V$ such that $\mu(v)\not\in V$).
%
%
%
%
%
%
%
\item[{\em Case 2}:] {\em Vertices of some copy of $G_i$
that are mapped by $\mu'$ to a single copy of $G_i$,
while not preserving their indices inside $G_i$.}

(This refers to vertices of some $G_i^{u,v}$
that are mapped by $\mu'$ to vertices of $G_i^{u',v'}$,
where $(u',v')$ may but need not equal $(u,v)$,
such that for some $j\in[k]$ the $j^\xth$ vertex of $G_i^{u,v}$
is not mapped by $\mu'$ to the $j^\xth$ vertex of $G_i^{u',v'}$.)
\footnote{Recall that $G_i^{u,v}$ and $G_i^{u',v'}$ are both
copies of the $k$-vertex graph $G_i$, which is an asymmetric graph,
and so the notion of the $j^\xth$ vertex in them is well-defined.
Formally, the $j^\xth$ vertex of $G_i^{u,v}$ is $\phi^{-1}(j)$ such
that $\phi$ is the (unique) bijection satisfying $\phi(G_i^{u,v})=G_i$.}

By the fact that $G_i$ is asymmetric,
it follows that each such copy contributes at least one unit
to the difference between~$G'$ and $\mu'(G')$,
and so (again) the total contribution of all these vertices
is proportional to their number; that is,
if the number of vertices in this case is $n_2$,
then their contribution is at least $n_2/k$.

%
%
\item[{\em Case 3}:]
{\em Vertices $v\in V$ such that $\mu'(v)\in V\setminus\{v\}$.}

(This is the main case, and here we use the hypothesis that
the edge-colored multi-graph~$G$ is robustly self-ordered.

%
Intuitively,
the hypothesis that the edge-colored~$G$ is robustly self-ordered
implies that such vertices contribute proportionally
to the difference between the colored versions of
the multi-graphs~$G$ and $\mu(G)$,
where $\mu$ is the restriction of $\mu'$ to $V$.
Indeed, we first assume, for simplicity, that $\mu'(V)=V$,
an assumption we shall have to dispose of later.
In this case, the number of tuples $(\{u,w\},i)$
such that $\{u,w\}$ is colored $i$ in exactly one of these multi-graph
(i.e., either in~$G$ or in $\mu(G)$ but not in both)
is at least $\gamma\cdot|\{v\!\in\!V:\mu(v)\neq v\}|$.
Assuming, without loss of generality that $\chi(\{u,w\})=i$
but either $\{\imu(u),\imu(w)\}\not\in E$
or $\chi(\{\imu(u),\imu(w)\})=j\neq i$,
we observe that $\imu(u)$ and $\imu(w)$ cannot
be connected in $G'$ via a copy of $G_i$.
We consider two sub-cases:
\BE
\item $\mu'$ maps a copy of $G_i$ to $G_i^{u,w}$,
but either $\imu(u)$ or $\imu(w)$ is not connected to this copy in $G'$.
In this sub-case we get a contribution of at least one unit,
since $u$ and $w$ are connected to $G_i^{u,w}$ in $G'$.
\item $\mu'$ does not map a copy of $G_i$ to $G_i^{u,w}$.
In this sub-case, it follows that some vertices that do not belong
to a copy of $G_i$ are mapped by $\mu'$ to $G_i^{u,w}$
which means that Case~1 applies for each such a tuple.
\EE
Hence, if the number of vertices in the current case is $n_3$,
then the number of tuples (handled by the two sub-cases)
is at least $\gamma\cdot n_3$, and we get a contribution
of at least $\gamma\cdot{n_3}/{k}$
(since the second sub-case is handled via Case~1).

The foregoing description is based on the assumption that $\mu(V)=V$.
If this does not hold, then we redefine $\mu$ such that $\mu(v)\not\in V$
is modified such that $\mu(v)=r$ if $r\in V$ has no preimage under $\mu'$.
(Of course, each such $r$ is only used once.)
Indeed, the modified $\mu$ may be ficticiously charged with $d$ edges
per each modification, but each such modification arises due to $v\in V$
that contributes at least one unit via (the last part of) Case~1.
Hence, the amortized over-counting of $d\cdot \gamma/k$ units
is compansated by the unit contributed in Case~1.
%

\item[{\em Case 4}:] {\em Vertices of some copy of $G_i$
that are mapped by $\mu'$ to a different copy of $G_i$.}

This refers to the case that $\mu'$ maps $G_i^{u,v}$
to $G_i^{u',v'}$ such that $(u',v')\neq(u,v)$,
which corresponds to mapping the gadget to a gadget
connecting a different pair of vertices
(but by an edge of the same color).

For $u,v,u',v'$ and $i$ as above, if $\mu'(u)=u'$ and $\mu'(v)=v'$,
then a gadget that connects $u$ and~$v$ in $G'$
is mapped to a gadget that does not connects them in $\mu'(G')$
(but rather connects the vertices $u'$ and $v'$,
whereas either $u'\neq u$ or $v'\neq v$).
So, due to the gadget-edge incident at either $u$ or $v$,
we get a contribution of at least one unit to the
difference between~$G'$ and $\mu'(G')$,
whereas the number of vertices in this gadget is $k$.
Hence, the contribution is proportional to the
number of non-fixed-points of the current type.
Otherwise (i.e., $(\mu'(u),\mu'(v))\neq(u',v')$),
we get a vertex as in Case~3,
and get a proportional contribution again.
%
%
%
%
%
%
%
%
%
%
\EDes
Hence, the contribution of each of these cases
to the difference between~$G'$ and $\mu'(G')$
is proportional to the number of vertices involved.
Specifically, if there are $n_i$ vertices in Case~$i$,
then we get a contribution-count of
at least $\gamma\cdot\sum_{i\in[4]}n_i/k$,
%
where some of these contributions were possibly counted thrice.
The claim follows.
\EPF

\BR[Fitting any desired number of vertices]
\label{colored2std:rem}
Assuming that the hypothesis of Theorem~\ref{colored2std:clm}
can be met for any sufficiently large $n\in S\subseteq\N$,
Construction~\ref{colored2standard:ct} yields robustly
self-ordered $n'$-vertex graphs for any $n'\in\{k\cdot n:n\!\in\!S\}$,
where $k=k(d,c)$ is as in Theorem~\ref{colored2std:clm}.
To obtain such graphs also for $n'$ that is not a multiple of $k$,
we may use two gadgets with a different number of vertices
for replacing at least one of the sets of colored edges.
\ER

\subsection{Application: Making the graph regular and expanding}
\label{edge-colored:applic:sec}
We view the edge-colored model as an intermediate locus
in a two-step methodology for constructing
robustly self-ordered graphs of bounded-degree.
First, one constructs edge-colored multi-graphs that
are robustly self-ordered
in the sense of Definition~\ref{colored-robust:def},
and then converts them to ordinary robustly self-ordered graphs
(in the sense of Definition~\ref{robust-asymmetric:def}),
by using Construction~\ref{colored2standard:ct}
(while relying on Theorem~\ref{colored2std:clm}).

We demonstrate the usefulness of this methodology
by showing that it yields a simple way of
making robustly self-ordered graphs be also expanding
as well as regular, while maintaining a bounded degree.
%
We just augment the original graph by super-imposing
an expander (on the same vertex set),
while using one color for the edges of the original graph
and another color for the edges of the expander.
Note that we do not have to worry about the possibility of
creating parallel edges (since they are assigned different colors).
The same method applies in order to make the graph regular.
We combine both transformations in the following result,
which we shall use in the sequel.

\BT[Making the graph regular and expanding]
\label{make:regular+expanding:thm}
For constant $d\geq3$ and~$\gamma$,
there exists an efficient algorithm
that given a $\gamma$-robustly self-ordered
graph $G=(V,E)$ of maximum degree $d$,
returns a $(d+O(1))$-regular multi-graph expander
coupled with a 2-coloring of its edges
such that the edge-colored multi-graph is $\gamma$-robustly self-ordered
{\rm(in the sense of Definition~\ref{colored-robust:def})}.
\ET
The same idea can be applied to edge-colored multi-graphs;
in this case, we use one color more than given.
We could have avoided the creation of parallel edges
with the same color by using more colors,
but preferred to relegate this task to
Construction~\ref{pre-colored2standard:ct},
while recalling that it preserves both the expansion
and the degree-regularity.
Either way, applying Theorem~\ref{colored2std:clm}
to the resulting edge-colored multi-graph,
we obtain robustly self-ordered (uncolored) graphs.
\medskip


\BPF
For any $d''\geq d+d'$,
given a graph $G=(V,E)$ of maximum degree $d$ that
is $\gamma$-robustly self-ordered and a $d'$-regular expander
graph $G'=(V,E')$,
%
we construct the desired $d''$-regular multi-graph $G''$
by super-imposing the two graphs on the same vertex set,
while assigning the edges of each of these graphs a different color.
In addition, we add edges to make the graph regular,
and color them using the same color as used for the expander.
\footnote{We assume for simplicity that $|V'|$ is even.
Alternatively, assuming that~$G$ contains no isolated vertex,
we first augment it with an isolated vertex
and apply the transformation on the resulting graph.
Yet another alternative is to consider only even $d''$.}
%
%
%
%
%
Details follow.

\BI
\item
We superimpose~$G$ and $G'$ (i.e., create a multi-graph $(V,E\cup E')$),
while coloring the edges of~$G$ (resp., $G'$)
with color~1 (resp., color~2).

Note that this may create parallel edges,
but with different colors.
\item
Let $d_v\leq d+d'$ denote the degree of vertex $v$
in the resulting multi-graph.
Then, we add edges to this multi-graph so that each vertex has degree~$d''$.
These edges will also be colored~2.

(Here, unless we are a bit careful, we may introduce
parallel edges that are assigned the same color.
This can be avoided by using more colors for these added edges,
%
%
but in light of Construction~\ref{pre-colored2standard:ct}
(which does essentially the same)
there is no reason to worry about this aspect.)
\EI
(Recall that the resulting edge-colored multi-graph is denoted $G''$.)

The crucial observation is that, since the edges of~$G$
are given a distinct color in $G''$, the added edges
do not harm the robust self-ordering feature of~$G$.
Hence, for any permutation $\mu:V\to V$,
any vertex-pair that contributes to the symmetric difference
between~$G$ and $\mu(G)$, also contributes to an inequality
between colored edges of $G''$ and $\mu(G'')$
(by virtue of the edges colored~1).
\EPF


\subsection{Local computability of the transformations}
\label{edge-colored:local:sec}
In this subsection, we merely point out that the transformation
presented in Constructions~\ref{pre-colored2standard:ct}
and~\ref{colored2standard:ct} as well as the one underlying
the proof of Theorem~\ref{make:regular+expanding:thm}
preserve efficient local computability
(e.g., one can determine the neighborhood of a vertex in
the resulting multi-graph by making a polylogarithmic number
of neighbor-queries to the original multi-graph).
Actually, this holds provided that we augment the
(local) representation of graphs, in a natural manner.

Recall that the standard representation of bounded-degree graphs
is by their incidence functions.
Specifically, a graph $G=([n],E)$ of maximum degree $d$ is represented
by the {\sf incident function} $g:[n]\times[d]\to[n]\cup\{0\}$
such that $g(v,i)=u\in[n]$ if $u$ is the $i^\xth$ neighbor of $v$,
and $g(v,i)=0$ if $v$ has less than $i$ neighbors.
This does not allow us to determined the identity of the $j^\xth$
edge in~$G$, nor even to determine the number of edges in~$G$,
by making a polylogarithmic number of queries to $g$,
whereas this determination is needed for
a local implementation of Construction~\ref{colored2standard:ct}.
Nevertheless, efficient local computability is preserved
if we use the following local representation
(presented for edge-colored multi-graphs).

%
\BD[Local representation]
For $d,c\in\N$, a {\sf local representation}
of a multi-graph $G=([n],E)$ of maximum degree $d$
that is coupled with a coloring $\chi\!:\!E\!\to\![c]$
is provided by the following three functions:
\BE
\item
An incidence function $g_1:[n]\times[d]\to\N\cup\{0\}$
such that $g_1(v,i)=j\in\N$
if $j$ is the index of the $i^\xth$ edge that incident at vertex $v$,
and $g_1(v,i)=0$ if $v$ has less than $i$ incident edges.
\item
An edge enumeration function $g_2:\N\to(\binom{[n]}{2}\times[c])\cup\{0\}$
such that $g_2(j)=(\{u,v\},\chi(e_j))$ if the $j^\xth$ edge,
denoted $e_j$, connects the vertices $u$ and~$v$,
and $g_2(j)=0$ if the multi-graph has less than $j$ edges.
\item
An vertex enumeration {\rm(by degree)}
function $g_3:[d]\times[n]\to[n]\cup\{0\}$
such that $g_3(i,j)=v\in[n]$
if $v$ is the $j^\xth$ vertex of degree $i$ in the multi-graph,
and $g_3(i,j)=0$ if the multi-graph has less than $j$ vertices of degree $i$.
\EE
\ED

The aforementioned incident function $g:[n]\times[d]\to[n]\cup\{0\}$
can be computed by composing $g_1$ and $g_2$;
in particular, $g(v,i)=u\in[n]$ if $g_2(g_1(v,i))=(\{u,v\},k)$
for some $k\in[c]$.
Needless to say, the function $g_3$ is redundant in the case
that we are guaranteed that the multi-graph is regular.
%
%
One may augment the foregoing representation by providing also
the total number of edges, but this number can be determined
by binary search.
%

\BT[The foregoing transformations preserve local computability]
\label{local-computability:thm}
The local representation of the multi-graph that
result from Construction~\ref{pre-colored2standard:ct}
can be computed by making a polylogarithmic number of queries
to the given multi-graph.
The same holds for Construction~\ref{colored2standard:ct}
and for the transformation underlying
the proof of Theorem~\ref{make:regular+expanding:thm}.
\ET

\BPF
For Construction~\ref{pre-colored2standard:ct},
we mostly need to enumerate all parallel edges
that connect~$u$ and~$v$. This can be done easily
by querying the incidence function on $(u,1),\ldots,(u,d)$
and querying the edge enumeration function on the non-zero answers.
(In addition, when adding a self-loop on vertex $v\in[n]$,
we need to determine the degree of $v$
as well as the number of edges in the multi-graph
(in order to know how to index the self-loop
in the incidence and edge enumeration functions, respectively).)

For Construction~\ref{colored2standard:ct},
we merely need to determine the color of the $j^\xth$ edge,
its endpoint
(and its index in the incidence list of each of its endpoints),
in order to replace this colored edge by the relevant gadget.
Recall that the relevant gadget
uses the vertices $n+(j-1)\cdot k+1,\ldots,n+j\cdot k$
and its edges are determined by the color of the edge that it replaces.

For the transformation underlying
the proof of Theorem~\ref{make:regular+expanding:thm},
adding edges to make the multi-graph regular requires
determining the index of a vertex in the list of all vertices
of the same degree (in order to properly index the added edges).
Here is where we use the vertex enumeration (by degree) function.
(We also need a local procedure $I$ for transforming
a sorted $n$-long sequence $(d_1,\ldots,d_n)\in[d'']$
into an all-$d''$ sequence by making pairs of increments;
that is, given $j\in[D]$ such that $D=(d''n-\sum_{i\in[n]}d_i)/2$,
we should determine a pair $I(j)=(I_1(j),I_2(j))\in[n]^2$
such that for every $i\in[n]$ it holds
that $d_i+|I_1^{-1}(i)|+|I_2^{-1}(i)|=d''$.)
\EPF
























































































































































































































































































































































































































































































































































































































































































\setcounter{section}{3}\section{The Three-Step Construction}
\label{three-step:sec}
In this section we present a different construction
of bounded-degree graphs that are robustly self-ordered.
It uses totally different techniques than the ones utilized
in the construction presented in Section~3.
Furthermore, the current construction offers the flexibility
of obtaining either graphs that have small connected components
(i.e., of logarithmic size)
or graphs that are highly connected (i.e., are expanders).
Actually, one can obtain anything in-between
(i.e., $n$-vertex graphs that consist of $s(n)$-sized
connected components that are each an expander,
for any $s(n)=\Omega((\log n)/\log\log n)$).
We mention that robustly self-ordered bounded-degree graphs
with small connected components are used in the proof
of Theorem~5.2.

As stated in Section~1.1.2,
the current construction proceeds in three steps.
First, in Section~\ref{step1:sec}, we prove the existence
of robustly self-ordered bounded-degree graphs,
and observe that such $\ell$-vertex graphs can actually
be found in $\poly(\ell!)$-time (i.e., $\exp(\tildeO(\ell))$-time).
Next, setting $\ell=\Omega((\log n)/\log\log n)$,
we use these graphs as part of $2\ell$-vertex
connected components in an $n$-vertex
(robustly self-ordered bounded-degree)
graph that is constructed in $\poly(n)$-time
(see Section~\ref{step2:sec}).
Lastly, in Section~\ref{step3:sec},
we repeat this strategy using the graphs
constructed in Section~\ref{step2:sec},
and obtain exponentially larger graphs
that are locally constructible.

%
%
In addition, in Section~\ref{local-so:sec},
we show that the foregoing graphs can be locally self-ordered.
That is, given a vertex $v$ in any graph $G'=(V',E')$
that is isomorphic to the foregoing $n$-vertex graph
and oracle access to the incidence function of $G'$,
we can find in $\poly(\log n))$-time the vertex
to which this unique isomorphism maps $v$.
%

%
%
%
%
%
%
%

\subsection{Existence}
\label{step1:sec}
As stated above, we start with establishing the mere existence
of bounded-degree graphs that are robustly self-ordered.

\BT[Robustly self-ordered graphs exist]
\label{existence:thm}
For any sufficiently large constant~$d$,
there exists a family $\{G_n\}_{n\in\N}$
of robustly self-ordered $d$-regular graphs.
Furthermore, these graphs are expanders.
\ET
Actually, it turns out that random $d$-regular graphs
are robustly self-ordered; see Theorem~6.1.
Either way, given the existence of such $n$-vertex graphs,
they can actually be found in $\poly(n!)$-time,
by an exhaustive search. Specifically,
for each of the possible $o((dn)!)$ graphs,
we check the robust self-ordering condition
by checking all $n!-1$ relevant permutation.
(The expansion condition can be checked similarly,
by trying all $(0.5+o(1))\cdot2^n$ relevant subsets of~$[n]$.)

The proof of Theorem~\ref{existence:thm} utilizes
a simpler probabilistic argument than the one used
in the proof of Theorem~6.1.
This argument (captured by Claim~\ref{random-colored:clm})
refers to the auxiliary model of edge-colored multi-graphs
(see Definition~\ref{colored-robust:def})
and is combined with a transformation of
this model to the original model of uncolored graphs
(provided in Construction~\ref{colored2standard:ct}
and analyzed in Theorem~\ref{colored2std:clm}).
Indeed, the relative simplicity of Claim~\ref{random-colored:clm}
is mainly due to using the edge-colored model
(see digest at the end of Section~6).
\medskip

\BPF
To facilitate the proof, we present the construction
while referring to the edge-colored model presented in
Section~\ref{edge-colored:sec}.
We shall then apply Theorem~\ref{colored2std:clm}
and obtain a result for the original model
(of uncolored simple graphs).

For $m=n/O(1)$,
we shall consider $2m$-vertex multi-graphs
that consists of two $m$-vertex cycles,
using a different color for the edges of each cycle,
that are connected by $d'=O(1)$ random perfect matching,
which are also each assigned a different color.
(Hence, we use $2+d'$ colors in total.)
We shall show that (w.h.p.) a random multi-graph constructed in this way
is robustly self-ordered (in the colored sense).
(Note that parallel edges, if they exist,
will be assigned different colors.)
%
Specifically,
we consider a generic $2m$-vertex multi-graph that is determined
by $d'$ perfect matchings of $[m]$ with $\{m+1,\ldots,2m\}$.
Denoting this sequence of perfect matchings by ${\ov M}=(M_1,\ldots,M_{d'})$,
we consider the (edge-colored) multi-graph $G_{\ov M}([2m],E_{\ov M})$
given by
\begin{eqnarray*}
E_{\ov M} &=& C_1\cup C_2 \cup\bigcup_{j\in[d']} M_j \\
& &\mbox{\rm where $C_1=\{\{i,i+1\}:i\in[m-1]\}\;\cup\;\{\{m,1\}\}$} \\
& &\mbox{\rm and $C_2=\{\{m+i,m+i+1\}:i\in[m-1]\}\;\cup\;\{\{2m,m+1\}\}$}
\end{eqnarray*}
and a coloring $\chi$ in which the edges of $C_j$ are colored $j$
and the edges of $M_j$ are colored $j+2$.
(That is, for $i\in\{1,2\}$, the set $C_i$ forms a cycle
of the form $((i-1)m+1,(i-1)m+2,\ldots,\allowbreak(i-1)m+m,(i-1)m+1)$
and its edges are colored~$i$.)
Note that the $d'+1$ edges incident at each vertex are assigned $d'+1$
different colors.

\Bcm{\em(W.h.p., $G_{\ov M}$ is robustly self-ordered):}
\label{random-colored:clm}
For some constant $\gamma>0$,
with high probability over the choice of~$\ov M$,
the edge-colored multi-graph $G_{\ov M}$ is $\gamma$-robustly self-ordered.
Furthermore, it is also an expander.
\Ecm

\Bpf
%
%
Consider an arbitrary permutation $\mu:[2m]\to[2m]$,
and let $t=|\{i\!\in\![2m]\!:\!\mu(i)\neq i\}|$.
%
%
We shall show that, with probability $1-\exp(-\Omega(dt\log m))$
over the choice of $\ov M$, the difference between the colored versions
of $G_{\ov M}$ and $\mu(G_{\ov M})$ is $\Omega(t)$.
Towards this end, we consider two cases.
\BDes
\item[{\em Case 1}:] $|\{i\in[m]:\mu(i)\not\in[m]\}|>t/4$.
Equivalently, $|\{i\in[2m]:\ceil{\mu(i)/m}\neq\ceil{i/m}\}|>t/2$.

The vertices in the set $\{i\in[m]:\mu(i)\not\in[m]\}$
are mapped from the first cycle to the second cycle,
and so rather than having two incident edges
that are colored~1 they have two incident edges colored~2.
%
Hence, each such vertex contributes two units to the difference
(between the colored versions of $G_{\ov M}$ and $\mu(G_{\ov M})$),
and the total contribution is greater than $2\cdot(t/4)\cdot 2$,
%
where the first factor of~2 accounts also for vertices
that are mapped from $C_2$ to $C_1$.

\item[{\em Case 2}:] $|\{i\in[m]:\mu(i)\not\in[m]\}|\leq t/4$.
Equivalently, $|\{i\in[2m]:\ceil{\mu(i)/m}\neq\ceil{i/m}\}|\leq t/2$.

We focus on the non-fixed-points of $\mu$ that stay on their
original cycle (i.e., those not considered in Case~1).
Let $A\eqdef\{i\!\in\![m]:\mu(i)\!\neq\!i\wedge\mu(i)\!\in\![m]\}$
and
$B\eqdef\{i\!\in\!\{m+1,....,2m\}:
   \mu(i)\!\neq\!i\wedge\mu(i)\!\in\!\{m+1,\ldots,2m\}\}$.
By the case hypothesis, $|A|+|B|\geq t/2$,
and we may assume (without loss of generality) that $|A|\geq t/4$.
As a warm-up, we first show that
{\em each element of~$A$ contributes
a non-zero number of units to the difference}\/
(between the colored versions of $G_{\ov M}$ and $\mu(G_{\ov M})$)
{\em with probability $1-O(1/m)^{d'}$}, over the choice of $\ov M$.

To see this, let $\pi_j:[m]\to\{m+1,\ldots,2m\}$
be the mapping used in the $j^\xth$ matching;
that is, $M_j=\{\{i,\pi_j(i)\}:i\!\in\![m]\}$,
which means that $\pi_j(i)$ is the $j^\xth$ match of $i$ in $G_{\ov M}$
(i.e., the vertex matched to $i$ by $M_j$).
Then, we consider the event that {\em for some $j\in[d']$,
the $j^\xth$ match of $i\in[m]$ in $\mu(G_{\ov M})$
is different from the $j^\xth$ match of $i$ in $G_{\ov M}$},
and note that when this event occurs $i$ contributes to the difference
(between the colored versions of $G_{\ov M}$ and $\mu(G_{\ov M})$).
Note that $x$ is the $j^\xth$ match of $i$ in $\mu(G_{\ov M})$
if and only if $\imu(x)$ is the $j^\xth$ match of $\imu(i)$ in $G_{\ov M}$,
which holds if and only if $\imu(x)=\pi_j(\imu(i))$
(equiv., $x=\mu(\pi_j(\imu(i)))$).
Hence, $i\in[m]$ contributes to the difference if and only if
for some $j$ it holds that $\pi_j(i)\neq\mu(\pi_j(\imu(i)))$,
because $\pi_j(i)\neq\mu(\pi_j(\imu(i)))$ means that
the edge $\{i,\pi_j(i)\}$ is colored~$j+2$ in $G_{\ov M}$
but is not colored~$j+2$ in $\mu(G_{\ov M})$
(since a different edge incident at~$i$ in $\mu(G_{\ov M})$
is colored~$j+2$).
%
%
Letting ${\ov\pi}=(\pi_1,\ldots,\pi_{d'})$,
the probability of the complementary event
(i.e., $i$ does not contribute to the difference)
is given by
\begin{eqnarray*}
\prob_{\ov\pi}
  \left[(\forall j\!\in\![d'])\;\;\pi_j(i)=\mu(\pi_j(\imu(i)))\right]
&=& \prod_{j\in[d']}
    \prob_{\pi_j}\left[\pi_j(i)=\mu(\pi_j(\imu(i)))\right] \\
&\leq& (m-1)^{-d'},
\end{eqnarray*}
where the inequality uses the hypothesis
that $\mu(i)\neq i$ and $i,\mu(i)\!\in\![m]$;
specifically, fixing the value of $\pi_j(\imu(i))$,
leaves $\pi_j(i)$ uniformly distributed
in $S\eqdef\{m+1,\ldots,2m\}\setminus\{\pi_j(\imu(i))\}$,
which means that
$\prob_{\pi_j}[\pi_j(i)\!=\!\mu(v)\,|\,v\!=\!\pi_j(\imu(i))]\leq1/|S|$
(where equality holds if $\mu(v)\in S$).

%
%
%
%
The same argument generalises to any set $I\subseteq A$
such that $I\cap\mu(I)=\emptyset$.
In such a case, letting $I=\{i_1,\ldots,i_{t'}\}$, we get
\begin{eqnarray*}
\lefteqn{\prob_{\ov\pi}\left[(\forall i\!\in\!I)(\forall j\!\in\![d'])\;\;
                     \pi_j(i)=\mu(\pi_j(\imu(i)))\right]} \\
&=& \prod_{k\in[t']}\prod_{j\in[d']}
    \prob_{\pi_j}\left[\pi_j(i_k)=\mu(\pi_j(\imu(i_k)))
          \left|(\forall k'\!\in\![k-1])\;
                   \pi_j(i_{k'})=\mu(\pi_j(\imu(i_{k'})))\right.\right] \\
&\leq& (m-2t'+1)^{-t'd'},
\end{eqnarray*}
where the inequality uses the hypothesis that $I\cap\mu(I)=\emptyset$;
specifically, for each $k\in[t']$, we use the fact
that $i_k\not\in\{i_1,\ldots,i_{k-1},\imu(i_1),\ldots,\imu(i_{k})\}$.
%
Hence, fixing the values of $\pi_j(i_{k'})$ for all $k'\in[k-1]$
and the values of $\pi_j(\imu(i_{k'}))$ for all $k'\in[k]$,
and denoting these values
by $u_1,\ldots,u_{k-1}$ and $v_1,\ldots,v_{k}$ respectively,
leaves $\pi_j(i_k)$ uniformly distributed
in $S\eqdef\{m+1,\ldots,2m\}\setminus\{u_1,\ldots,u_{k-1},v_1,\ldots,v_k\}$,
which means that
$\prob_{\pi_j}[\pi_j(i)\!=\!\mu(v_k)\,|\,\mbox{\rm foregoing fixing}]\leq1/|S|$
(where equality holds if $\mu(v_k)\in S$).
%
%
%
%
%
%
%

Recalling that $|A|\geq t/4$ and $t\leq2m$,
we upper-bound the probability (over the choice of $\ov M$)
that $A$ contains a $t/8$-subset $A'$ such that
$(\forall i\!\in\!A')(\forall j\!\in\![d'])\;\pi_j(i)=\mu(\pi_j(\imu(i)))$,
by taking a union bound over all possible $A'$
and using for each such $A'$ a subset $I\subset A'$
such that $I\cap\mu(I)=\emptyset$.
(So we actually take a union bound over the $I$'s
and derive a conclusion regarding the $t/8$-subsets $A'$.)
Observing that $|I|\geq|A'|/3\geq t/24$,
%
we conclude that, with probability at most
$\binom{t}{t/24}\cdot(m/2)^{d'\cdot t/24}=\exp(-\Omega(d' t\log m))$
over the choice of~$\ov M$,
the set $A$ contains no $t/8$-subset $A'$ as above.
This means that,
with probability at most $\exp(-\Omega(d' t\log m))$,
less than $t/8$ of the indices $i\in A$ contribute
a non-zero number of units to the difference
(between the colored versions of $G_{\ov M}$ and $\mu(G_{\ov M})$).
\EDes
Hence, we have shown that, for every permutations $\mu:[2m]\to[2m]$,
the probability (over the choice of $\ov M$) that the size
of the symmetric difference between
the colored versions of $G_{\ov M}$ and $\mu(G_{\ov M})$
is smaller than $t/8$ is $\exp(-\Omega(d' t\log m))$,
where $t$ is the number of non-fixed-points of $\mu$.
%
Letting $\gamma=1/8$ and taking a union bound over
all (non-trivial) permutations $\mu:[2m]\to[2m]$,
we conclude that the probability, over the choice of $\ov M$,
that $G_{\ov M}$ is not $\gamma$-robustly self-ordered is at most
\begin{eqnarray*}
\sum_{t\in[2m]}\binom{2m}{t}\cdot\exp(-\Omega(d' t\log m))
&=& \sum_{t\in[2m]}\exp(-\Omega((d'-O(1))\cdot t\log m)) \\
&=& \exp(-\Omega((d'-O(1))\cdot \log m)),
\end{eqnarray*}
and the claim follows (for any sufficiently large $d'$),
while observing that, with very high probability,
these multi-graphs are expanders.
\Epf

\mypar{Back to the non-colored version.}
We now convert the edge-colored multi-graphs $G=G_{\ov M}$
that are $\gamma$-robustly self-ordered into standard graphs $G'$
that are robustly self-ordered in the original sense.
This is done by using Construction~\ref{colored2standard:ct}
(while relying on Theorem~\ref{colored2std:clm}).
Recall that this transformation also preserves expansion.
Actually, before invoking Construction~\ref{colored2standard:ct},
we augment the multi-graph~$G$ by adding a self-loop to each vertex,
and color all these self-loops using a special color.
Combining Claim~\ref{random-colored:clm}
and Theorem~\ref{colored2std:clm}, the current theorem follows.
%
\EPF

\subsection{Constructions}
\label{step2:sec}
Having established the existence of bounded-degree graphs that
are robustly self-ordered, we now turn to actually construct them.
We shall use the fact that the proof of existence yields
a construction that runs in time that is polynomial in the
number of possible graphs.
Specifically, for $\ell=\frac{O(\log n)}{\log\log n}$,
we shall construct $\ell$-vertex graphs in $\poly(\ell^\ell)$-time
and use them in our construction of $n$-vertex graphs,
while noting that $\poly(\ell^\ell)=\poly(n)$.

\BT[Constructing robustly self-ordered graphs]
\label{construction:thm}  \label{CONSTRUCTION:THM}   
For any sufficiently large constant~$d$,
there exists an efficiently constructable family $\{G_n\}_{n\in\N}$
of robustly self-ordered graphs of maximum degree $d$.
That is, there exists a polynomial-time algorithm
that on input $1^n$ outputs the $n$-vertex graph $G_n=([n],E_n)$.
Furthermore, $G_n$ consists of connected components
of size $\frac{O(\log n)}{\log\log n}=o(\log n)$.
%
\ET
Note that the connected components of $G_n$ cannot be any smaller
(than $\frac{O(\log n)}{\log\log n}$).
This is the case because an asymmetric $n$-vertex bounded-degree graph,
let alone a robustly self-ordered one, cannot have connected components
of size $\frac{o(\log n)}{\log\log n}$ (because the number of $t$-vertex
graphs of bounded-degree is $t^{O(t)}$).
%
\medskip

\BPF
The proof proceeds in two steps.
We first use the existence of $\ell$-vertex ($d'$-regular)
expander graphs that are robustly self-ordered
towards constructing a sequence of $m=\exp(\Omega(\ell\log\ell))$
bounded-degree $2\ell$-vertex graphs that are robustly self-ordered,
expanding, and far from being isomorphic to one another.
We construct this sequence of $2\ell$-vertex graphs in $\poly(m)$-time,
using the fact that $(\ell!)^{O(1)}=\poly(m)$.
In the second step, we show that the $(m\cdot2\ell)$-vertex graph
that consists of these $2\ell$-vertex graphs (as its connected components)
is robustly self-ordered.
Note that this graph is constructed in time that is
polynomial in its size, since its size is $\Omega(m)$,
whereas it is constructed in $\poly(m)$-time.
\footnote{We mention that a slightly different construction can be
based on the fact that random $\ell$-vertex ($d'$-regular) graphs
are robustly self-ordered expanders (see Theorem~6.1).
In this alternative construction we find a sequence of $m$ such
graphs that are pairwise far from being isomorphic to one another.
As further detailed in Remark~6.2,
the analysis of the alternative construction is somewhat easier
than the analysis of the construction presented below,
but we need the current construction for the proof
of Theorem~\ref{strong-construction:thm}.}

Given a generic $n$, let $\ell=\frac{O(\log n)}{\log\log n}$,
which implies that $\ell^\ell=\poly(n)$.
By Theorem~\ref{existence:thm}, for all sufficiently large $d'$,
there exist $\ell$-vertex $d'$-regular expander graphs
that are robustly self-ordered
(with respect to the robustness parameter $c'$).
Furthermore, we can find such a graph, denoted $G'_\ell$,
in time $\poly(\ell^\ell)=\poly(n)$,
by scanning all $\ell$-vertex $d'$-regular graphs and checking
both the expansion and the robustness (w.r.t parameter $c'$)
conditions for each of them.
%
%
%
Actually, for $d''=d'+1$, we shall also find
an $\ell$-vertex $d''$-regular expander, denoted $G''_\ell$,
that is robustly self-ordered.

\mypar{The construction of $G_n$.}
Using $G'_\ell$ and $G''_\ell$, we construct an $n$-vertex
robustly self-ordered graph, denoted $G_n$,
that consists of $n/2\ell$ connected components that
are pairwise far from being isomorphic to one another.
This is done by picking $m=n/2\ell$ permutations,
denoted $\pi_1,\ldots,\pi_m:[\ell]\to[\ell]$,
that are pairwise far-apart
and constructing $2\ell$-vertex graphs
such that the $i^\xth$ such graph consist of
a copy of $G'_\ell$ and a copy of $G''_\ell$ that are connected
by a matching as determined by the permutation $\pi_i$.
Specifically,
for $G'_\ell=([\ell],E'_\ell)$ and $G''_\ell=([\ell],E''_\ell)$,
the $i^\xth$ connected component is isomorphic
to a graph with the vertex set $[2\ell]$
and the edge set
\begin{equation}\label{edges-of-cc:eq1}
E'_\ell \;\cup\; \{\{\ell+u,\ell+v\}:\{u,v\}\in E''_\ell\}
         \;\cup\; \{\{v,\ell+\pi_i(v)\}:v\!\in\![\ell]\}.
\end{equation}
(The first two sets correspond to the copies of~$G'_\ell$ and $G''_\ell$,
and the third set corresponds to the matching between these copies.
Note that the vertices in $[\ell]$ have degree $d'+1$,
whereas vertices in $\{\ell+1,\ldots,2\ell\}$ have degree $d''+1\neq d'+1$.)


To see that this construction can be carried out in $\poly(n)$-time,
we need to show that the sequence of $m$ pairwise far-apart permutations
can be determined in $\poly(n)$-time, let alone that such a sequence exists.
This is the case, because we can pick the permutation sequentially
(one after the other) by scanning the symmetric group on $[\ell]$
and relying on the fact that for ($i<n$ and) any fixed sequence
of permutations $\pi_1,\ldots,\pi_{i-1}:[\ell]\to[\ell]$
it holds that a random permutation $\pi_i$ is far-apart from
each of the fixed $i-1$ permutations; that is,
$\prob_{\pi_i}[|\{v\in[\ell]:\pi_i(v)\neq\pi_j(v)\}|=\Omega(\ell)]=1-o(1/n)$
for every $j\in[i-1]$.
\footnote{Specifically, for some $\ell'=\Omega(\ell)$,
we upper-bound $\prob_{\pi}[|\{v\in[\ell]:\pi(v)=v)\}|\geq\ell-\ell']$,
where $\pi:[\ell]\to[\ell]$ is a random permutation.
We do so by observing that the number of permutations
that have at least $\ell-\ell'$ fixed-points is at most
$\binom{\ell}{\ell'}\cdot(\ell'!)=\frac{\ell!}{(\ell-\ell')!}$,
whereas $(\ell-\ell')!=\exp(\Omega(\ell\log\ell))=\omega(n)$
for any $\ell'$ such that $\ell-\ell'=\Omega(\ell)$.}
%
%

%
%
%
%
%
%

\mypar{Towards proving that $G_n$ is robustly self-ordered.}
We now prove that the resulting graph~$G_n$,
which consists of these $m$ connected components,
is $c$-robustly self-ordered, where $c$ is a universal constant
(which is independent of the generic $n$).
For starters, let's verify that $G_n$ is self-ordered.
We first note that any automorphism of $G_n$ must
map the vertices of copies of $G'_\ell$ (resp., $G''_\ell$)
to vertices of copies of $G'_\ell$ (resp., $G''_\ell$),
since these are the only vertices of degree~$d'+1$ (resp., $d''+1$).
The connectivity of these copies implies that
the automorphism must map each connected component
to some connected component,
which  determines the $m$ connected components.
The self-ordered feature of $G'_\ell$ and $G''_\ell$
determines a unique ordering on each copy,
whereas the fact the permutations (i.e., $\pi_i$'s)
are different imposes that each connected
component is mapped to itself
(i.e., the order of the connected components is preserved).
Hence, the automorphism must be trivial
(and it follows that $G_n$ is self-ordered).

%
%
%
%
%
%
%
%
%
%
%

An analogous argument establishes the robust self-ordering of $G_n$,
where we use the hypothesis that $G'_\ell$ and $G''_\ell$ are expanders
(rather than merely connected),
the choice of the $\pi_i$'s as being far-apart
(rather than merely different),
and the robust self-ordering of $G'_\ell$ and $G''_\ell$
(rather than their mere self-ordering).
%
Considering an arbitrary permutation $\mu:[n]\to[n]$,
these stronger features are used to establish a lower bound on the size
of the symmetric difference between $G_n$ and $\mu(G_n)$ as follows:
\BI
\item
The fact that $G'_\ell$ is an expander implies that if $\mu$ splits
the vertices of a copy of $G'_\ell$ such that $\ell'$ vertices are mapped
to copies that are different than the other $\ell-\ell'\geq\ell'$ vertices,
then this contributes $\Omega(\ell')$ units
to the difference between $G_n$ and $\mu(G_n)$.
Ditto for $G''_\ell$, whereas mapping a copy of $G'_\ell$
to a copy of $G''_\ell$ contributes $\Omega(\ell)$ units
(per the difference in the degrees).
\item
The robust self-ordering of $G'_\ell$ and $G''_\ell$ implies that
if $\mu$ changes the index of vertices inside a component,
then this yields a proportional difference between $G_n$ and $\mu(G_n)$.
\item
The distance between the $\pi_i$'s
(along with the aforementioned robustness)
implies that if~$\mu$ changes the indices of the connected components,
then each such change contributes $\Omega(\ell)$ units
to the difference between $G_n$ and $\mu(G_n)$.
\EI
The actual implementation of this sketch requires
a careful accounting of the various contributions.
As a first step in this direction we provide
a more explicit description of $G_n$.
We denote the set of vertices
of the copy of $G'_\ell$ (resp., $G''_\ell$)
in the $i^\xth$ connected component of $G_n$
by $F_i=\{2(i-1)\ell+j:j\in[\ell]\}$
(resp., $S_i=\{2(i-1)\ell+\ell+j:j\in[\ell]\}$).
Recall that $F_i$ and $S_i$ are connected by the edge-set
\begin{equation}\label{edges-of-cc:eq2}
\left\{\{2(i-1)\ell+j,2(i-1)\ell+\ell+\pi_i(j)\}:j\!\in\![\ell]\right\}
\end{equation}
whereas the subgraph of $G_n$ induced by $F_i$ (resp., $S_i$)
has the edge-set $\{\{2(i-1)\ell+u,2(i-1)+v\}:\{u,v\}\!\in\!E'_\ell\}$
(resp., $\{\{2(i-1)\ell+\ell+u,2(i-1)+\ell+v\}:\{u,v\}\!\in\!E''_\ell\}$).
In addition, let $F=\bigcup_{i\in[m]}F_i$
(resp., $S=\bigcup_{i\in[m]}S_i$).

\mypar{The actual proof (that $G_n$ is robustly self-ordered).}
Considering an arbitrary permutation $\mu:[n]\to[n]$,
we lower-bound the distance
(i.e., size of the symmetric difference)
between $G_n$ and $\mu(G_n)$
as a function of the number of non-fixed-points under $\mu$
(i.e., the number of $v\in[n]$ such that $\mu(v)\neq v$).
We do so by considering the (average) contribution
of every non-fixed-point to the distance between $G_n$ and $\mu(G_n)$
(i.e., number of pairs of vertices that form an edge in one graph
but not in the other). We may include the same contribution in
few of the following (seven) cases, but this only means that
we are double-counting the contribution by a constant factor.
We first consider the set of vertices that switch sides
(under $\mu$) in the connected components.

\BDes
\item[{\em Case 1}:] {\em Vertices $v\in F$ such that $\imu(v)\in S$}.
Ditto for $v\in S$ such that $\imu(v)\in F$.

Each such vertex contributes at least one unit to the distance
(between $G_n$ and $\mu(G)$)
by virtue of $v$ having degree~$d'+1$ in $G_n$
and strictly higher degree in $\mu(G_n)$,
since vertices in $F$ have degree~$d'+1$ (in $G_n$)
whereas vertices in~$S$ have higher degree (in $G_n$).
\footnote{Note that $v$ neighbors $u$ in $\mu(G_n)$
if and only if $\imu(v)$ neighbors $\imu(u)$ in $G_n$.}
%
%
%
%
%
\EDes
In light of Case~1, we may focus on vertices whose ``type''
(or ``side'') is preserved by $\imu$. Actually, it will be more convenient
to consider the set of vertices whose ``type'' is preserved by $\mu$;
%
that is, the set
$\{v\!\in\!F\!:\!\mu(v)\!\in\!F\}\cup\{v\!\in\!S\!:\!\mu(v)\!\in\!S\}$.
%
Next, for each $i\in[m]$, we define $\si'(i)$ to be the index
of the connected component that takes the plurality of $\mu(F_i)$;
that is, $\si'(i)\eqdef j$
if $|\{v\in F_i:\mu(v)\in F_j\}|\geq|\{v\in F_i:\mu(v)\in F_k\}|$
for all $k\in[m]$ (breaking ties arbitrarily).
Next, we consider the set of vertices that are not mapped
according to the plurality vote.

\BDes
\item[{\em Case 2}:]
{\em Vertices $v\in F_i$ such that $\mu(v)\in F\setminus F_{\si'(i)}$}.

For starters, suppose that
$|\{v\!\in\!F_i\!:\!\mu(v)\!\in\!F_{\si'(i)}\}|\geq\ell/2$;
that is, a majority of the vertices of $F_i$ are mapped by~$\mu$
to $F_{\si'(i)}$. In this case, by the expansion of $G'_\ell$,
we get a contribution that is proportional to the size
of the set $F'_i\eqdef\{v\!\in\!F_i\!:\!\mu(v)\!\not\in\!F_{\si'(i)}\}$,
because (in $G_n$) there are $\Omega(|F'_i|)$ edges
between $F'_i$ and the rest of $F_i$
but there are no edges between $F'_i$ and $F_i\setminus F'_i$ in $\mu(G_n)$.
In the general case, we have to be more careful since expansion
is guaranteed only for sets that have size at most $\ell/2$.
In such a case we use an adequate subset of $F'_i$.
Details follow.


Let $R_i\eqdef\{v\!\in\!F_i\!:\!\mu(v)\!\in\!F\setminus F_{\si'(i)}\}$.
We first note that there exists a set $J\subseteq[m]\setminus\{\si'(i)\}$
such that $F'_i\eqdef\bigcup_{j\in J}\{v\!\in\!F_i\!:\!\mu(v)\!\in\!F_j\}$
satisfies $|R_i|/4\leq|F'_i|\leq\ell/2$.
%
%
%
%
%
%
%
%
%
%
%
%
%
%
%
%
%
%
%
Recall that the subgraph of $G_n$ induced by $F_i$ is an expander,
and consider the edges in $G_n$ that cross the cut
between $F'_i$ and the rest of $F_i$. 
Then, this cut has $\Omega(|F'_i|)$ edges in $G_n$,
but there are no edges between $F'_i$ and $F_i\setminus F'_i$ in $\mu(G_n)$,
because $\imu(F'_i)\subseteq\bigcup_{j\in J}F_j$
and $\imu(F_i\setminus F'_i)\subseteq\bigcup_{j\in[m]\setminus J}F_j$
are not connected in $G_n$.
%
%
%
%
%
%
Hence, the total contribution of the vertices
in $\{v\!\in\!F_i\!:\!\mu(v)\!\in\!F\setminus F_{\si'(i)}\}=R_i$
to the distance (between $G_n$ and $\mu(G)$) is $\Omega(|F'_i|)$,
which is proportional to their number
(i.e., it is $\Omega(|R_i|)$).
%
\EDes
Defining $\si''(i)$ in an analogous manner with respect to $\mu(S_i)$,
we get an analogous contribution by the expander induced by $S_i$.
Specifically, for each $i\in[m]$, we define $\si''(i)$ to be the index
of the connected component that takes the plurality of $\mu(S_i)$;
that is, $\si''(i)\eqdef j$ if
$|\{v\!\in\!S_i\!:\!\mu(v)\!\in\!S_j\}|
  \geq|\{v\!\in\!S_i\!:\!\mu(v)\!\in\!S_k\}|$
for all $k\in[m]$ (breaking ties arbitrarily).
Analogously to Case~2, we consider the following vertices.

\BDes
\item[{\em Case 3}:]
{\em Vertices $v\in S_i$ such that $\mu(v)\in S\setminus S_{\si''(i)}$}.

Here we get a contribution of
$\Omega(|\{v\!\in\!S_i\!:\!\mu(v)\!\in\!S\setminus S_{\si''(i)}\}|)$,
where the analysis is analogous to Case~2.
\EDes
Recall that if $v\in F_i$ then it holds that $v=2(i-1)\ell+j$
for some $j\in[\ell]$, and that (in $G_n$) vertex $v$ has
a unique neighbor in~$S$, which is $2(i-1)\ell+\ell+\pi_i(j)\in S_i$.
It will be convinient to denote this neighbor by $\phi_i(v)$;
that is, for $v\in F_i$ such that $v=2(i-1)\ell+j$,
we have $\phi_i(v)=2(i-1)\ell+\ell+\pi_i(j)\in S_i$.
The next two cases refer to vertices that are mapped by $\mu$
according to the plurality vote
(e.g., $v\in F_i$ is mapped to $\mu(v)\in F_{\si'(i)}$),
but their match is not mapped accordingly
(i.e., $\phi_i(v)\in S_i$ is not mapped to $S_{\si'(i)}$).

\BDes
\item[{\em Case 4}:]
{\em Vertices $v\in F_i$ such that $\mu(v)\in F_{\si'(i)}$
but $\mu(\phi_i(v))\not\in S_{\si'(i)}$}.
%
%
%

(Note that the condition $v\in F_i$ and $\mu(v)\in F_{\pi'(i)}$
means that vertex $v$ is not covered in Case~2.
If $\si''(i)=\si'(i)$, then $\mu(\phi_i(v))\not\in S_{\si'(i)}$
means that $v$ is covered in Case~3, since $\phi_i(v)\in S_{i}$.
Hence, the current case is of interest only when $\si''(i)\neq\si'(i)$.
In particular, it is of interest when referring to vertices
in the $i^\xth$ connected component of $G_n$ that reside
in the copies of~$G'_\ell$ and $G''_\ell$ and are mapped
according to the plurality votes of these copies,
whereas these two plurality votes are inconsistent.)

We focus on the case that a vast majority of the vertices in
both $F_i$ and $S_i$ are mapped according to the plurality votes
(i.e., $\si'(i)$ and $\si''(i)$), since the complementary cases
are covered by Cases~2 and~3, respectively.
Specifically, if
either $|\{v\!\in\!F_i\!:\!\mu(v)\!\in\![n]\setminus F_{\si'(i)}\}|>\ell/3$
or $|\{u\!\in\!S_i\!:\!\mu(u)\!\in\![n]\setminus S_{\si''(i)}\}|>\ell/3$,
then we get a contribution of $\Omega(\ell)$
either by Cases~1\&2 or by Cases~1\&3.
%
Otherwise, it follows that
$${|\{v\!\in\!F_i\!:\!\mu(v)\!\in\!F_{\si'(i)}
    \;\wedge\; \mu(\phi_i(v))\!\in\!S_{\si''(i)}\}|}
    \geq \ell -2\cdot\ell/3$$
which implies that, if $\si'(i)\neq\si''(i)$,
then the $i^\xth$ connected component of $G_n$
contributes $\ell/3$ units to the difference
(between $G_n$ and $\mu(G_n)$),
since $v$ and $\phi_i(v)$ are connected in $G_n$,
but $\mu(v)\in F_{\si'(i)}$ and $\mu(\phi_i(v))\in S_{\si''(i)}$
reside in different connected components of $\mu(G_n)$.
(That is, the contribution is due to vertices $v$ of $F_i$
that are mapped by $\mu$ to $F_{\si'(i)}$,
while the corresponding vertices $\phi_i(v)$ of $S_i$
(which are connected to them in $G_n$) are mapped by $\mu$
to $S_{\si''(i)}\subset S\setminus S_{\si'(i)}$,
whereas $F_{\si'(i)}$ and $S_{\si''(i)}$ are not connected in $G_n$,
assuming $\si'(i)\neq\si''(i)$.)
%
%
%
%


To conclude: The contribution of the vertices of Case~4
(to the difference between $G_n$ and $\mu(G_n)$)
is proportional to the number of these vertices
(where this contribution {\em might}\/ have been counted
already in Cases~1,~2 and~3).

\item[{\em Case 5}:]
{\em Vertices $v\in F_i$ such that $\mu(v)\not\in F_{\si''(i)}$
but $\mu(\phi_i(v))\in S_{\si''(i)}$}.
%
%
%

(Equiv., vertices $v\in S_i$ such that $\mu(v)\in S_{\si''(i)}$
but $\mu(\phi_i^{-1}(v))\not\in F_{\si''(i)}$.)

Analogously to Case~4, the contribution of these vertices
is proportional to their number.
(Analogously, this augments Case~2 only in case $\si''(i)\neq\si'(i)$.)
\EDes
In light of Cases~2--5, we may focus on indices $i\in[m]$
such that $\si'(i)=\si''(i)$ and on vertices in the $i^\xth$
connected component that are mapped by $\mu$ to
the $\si'(i)^\xth$ connected component
(and the same ''type'' per Case~1).
The following case refers to such vertices that do not
maintain their position in this connected component.

\BDes
\item[{\em Case 6}:]
{\em Vertices $v\!=\!2(i-1)\ell+j\in F_i$
such that $\mu(v)\in F_{\si'(i)}\setminus\{2(\si'(i)-1)\ell+j\}$.}

Ditto for $v\!=\!2(i-1)\ell+\ell+j\in S_i$
such that $\mu(v)\in S_{\si''(i)}\setminus\{2(\si''(i)-1)\ell+\ell+j\}$.

(This case refers to vertices in $F_i$ that are mapped to $F_{\si'(i)}$
but do not maintain their index in the relevant copy of $G'_\ell$;
indeed, $v\!=\!2(i-1)\ell+j$ is the $j^\xth$ vertex of $F_i$,
but it is mapped by~$\mu$ to the $k^\xth$ vertex of $F_{\si'(i)}$
(i.e., $\mu(v)\!=\!2(\si'(i)-1)\ell+k$) such that $k\neq j$.)
%
%

Fixing $i$, let
$C\eqdef\{v\!=\!2(i-1)\ell+j\in F_i:
  \mu(v)\in F_{\si'(i)}\setminus\{2(\si'(i)-1)\ell+j\}\}$
denote the set of vertices considered in this case,
and $D=\{v\in F_i:\mu(v)\not\in F_{\si'(i)}\}$
denote the set of vertices that we are going to discount for.
%
As a warm-up, consider first the case that $D=\emptyset$.
In this case, by the robust self-ordering of $G'_\ell$,
the contribution of the vertices in $C$
to the difference between $G_n$ and $\mu(G_n)$
is $\Omega(|C|)$.

In the general case (i.e., where $D$ may not be empty),
we get a contribution of $\Omega(|C|)-d'\cdot|D|$,
where the second term compensates for the fact that the vertices
of $D$ were moved outside of this copy of $G'_\ell$ and replaced
by different vertices that may have different incidences.
Letting~$c'$ be the constant hidden in the $\Omega$-notation,
we get a contribution of
at least $c'\cdot|C|-d'\cdot|D|$,
which is at least $c'\cdot|C|/2$ if $|D|\leq c'\cdot|C|/2d'$.
On the other hand, if $|D|>c'\cdot|C|/2d'$,
then we get a contribution of $\Omega(|D|)=\Omega(|C|)$ by Cases~1--2.

Hence, in both sub-cases we have a contribution of $\Omega(|C|)$
to the difference between $G_n$ and $\mu(G_n)$.

The same analysis applies to
$\{v\!=\!2(i-1)\ell+\ell+j\in S_i:
  \mu(v)\in S_{\si''(i)}\setminus\{2(\si''(i)-1)\ell+\ell+j\}\}$,
where we use the robust self-ordering of $G''_\ell$
and Cases~1\&3.
%
%
\EDes
Lastly, we consider vertices that do not fall into any of
the prior cases. Such vertices maintain their type,
are mapped with the plurality vote of their connected component,
which is consistent among its two parts (i.e., $\si'$ and $\si''$),
and maintain their position in that component.
Hence, {\em the hypothesis that they are not fixed-points of $\mu$
can only be attributed to the fact that these vertices are
mapped to a connected component with a different index}.

\BDes
\item[{\em Case 7}:]
{\em Vertices $v\in F_i$ such that
both $\mu(v)\in F_{\si'(i)}\setminus F_{i}$
and $\mu(\phi_i(v))\in S_{\si''(i)}\setminus S_{i}$ hold.}
%

(We may assume that $\si'(i)\neq i$ and $\si''(i)\neq i$,
since otherwise this set is empty.
%
%
%
%
We may also assume that $\si'(i)=\si''(i)$,
since the complementary case was covered by Cases~4 and~5.
Hence, we focus on pairs of vertices that are
matched in the $i^\xth$ connected component of $G_n$
and are mapped by $\mu$ to the $k^\xth$ component of $G_n$
such that $k\neq i$.)

For every $i\neq k$,
let $\Delta_{i,k}=\{j\in[\ell]:\pi_i(j)\neq\pi_{k}(j)\}$
be the sets on which $\pi_i$ and $\pi_k$ differ.
(Note that if for every $v=2(i-1)\ell+j\in F_i$
it holds that $\mu(v)=2(k-1)\ell+j$
and $\mu(\phi_i(v))=2(k-1)\ell+\pi_i(j)$
(equiv., $\mu(2(i-1)\ell+\ell+\pi_i(j))=2(k-1)\ell+\pi_i(j)$),
then we get a contribution of $|\Delta_{i,k}|$
to the difference between $G_n$ and $\mu(G_n)$.)

%
Fixing $i$, let $D=D_1\cup D_2$ such that
\begin{eqnarray*}
%
%
%
%
%
%
%
D_1&=& \{v\in F_i:\mu(v)\not\in F_{\si'(i)}\;\vee\;
                 \mu(v+\ell)\not\in S_{\si''(i)}\} \\[3pt]
D_2&=& \left\{v\!=\!2(i-1)\ell+j\in F_i:
        \begin{array}{l}
         \mu(v)\in F_{\si'(i)}\setminus\{2(\si'(i)-1)\ell+j\} \\
         \vee\;\;
         \mu(\phi_i(v))
              \in S_{\si''(i)}\setminus\{2(\si''(i)-1)\ell+\ell+\pi_i(j)\}
        \end{array}\right\}
\end{eqnarray*}
(Recall that $\phi_i(2(i-1)\ell+j)=2(i-1)\ell+\ell+\pi_i(j)$.
The set $D_1$ accounts for the vertices covered in Cases~2\&3,
whereas $D_2$ accounts for the vertices covered in
(the two sub-cases of) Case~6.)

As a warm-up, consider first the case that $D=\emptyset$.
%
In this case, assuming $\si'(i)=\si''(i)\neq i$,
we get a contribution of $|\Delta_{i,\si'(i)}|=\Omega(\ell)$
(to the difference between $G_n$ and $\mu(G_n)$).
This contribution is due to the difference in the edges
that match $F_{\si'(i)}$ and $S_{\si'(i)}$ in $G_n$
and the edges that match $F_i$ and $S_i$ in $G_n$,
where $|\Delta_{i,\si'(i)}|=\Omega(\ell)$ is due to
the fact that the permutations (i.e., $\pi_k$'s) are far-apart.
The hypothesis $D_1=\emptyset$ means that all vertices
of $F_i$ (resp., of $S_i$) are mapped to $F_{\si'(i)}$
(resp., to $S_{\si''(i)}=S_{\si'(i)}$),
whereas $D_2=\emptyset$ means that these vertices preserves
their order within the two parts of the connected component.

The general case (i.e., where $D$ may not be empty)
requires a bit more care.
Suppose that the $\pi_k$'s are $\gamma$-apart;
that is, $|\Delta_{k',k}|>\gamma\cdot\ell$ for every $k'\neq k$.
We focus on the case that a vast majority of the vertices in
both $F_i$ and $S_i$ are mapped according to the plurality votes
(i.e., $\si'(i)$ and $\si''(i)$), since the complementary cases
are covered by Cases~2 and~3, respectively.
Specifically, if $|D_1|>\gamma\ell/3$,
then we get a contribution of $\Omega(\ell)$ by either Case~2 or Case~3.
Likewise, if $|D_2|>\gamma\ell/3$,
then we get a contribution of $\Omega(\ell)$ by Case~6.
%
So, assuming $\si'(i)=\si''(i)\neq i$, we are left with the case that
$$|\{v\!=\!2(i-1)\ell+j\in F_i\setminus D: j\in\Delta_{i,\si'(i)}\}|
  \geq\gamma\ell-2\gamma\ell/3.$$
In this case,
%
we get a contribution of at least $\gamma\ell/3$
to the difference between $G_n$ and $\mu(G_n)$.
This contribution is due to the difference in the edges
that match $F_{\si'(i)}$ and $S_{\si'(i)}$ in $G_n$
and the edges that match $F_i$ and $S_i$ in $G_n$,
where edges that have an endpoint (or its $\phi_i$-mate)
in $D$ were discarded.
Specifically, letting $k=\si'(i)=\si''(i)\neq i$,
the pair $(v,w)=(2(i-1)\ell+j,2(i-1)\ell+\ell+\pi_i(j))\in F_i\times S_i$
contributes to the difference if $j\in\Delta_{i,k}$ and
both $\mu(v)=2(k-1)\ell+j\in F_{k}$
and $\mu(w)=2(k-1)\ell+\ell+\pi_i(j)\in S_{k}$ hold
(i.e., $v\not\in D_1$ and $v,\phi_i^{-1}(w)\not\in D_2$).
\footnote{Recall
that $\phi_i^{-1}(w)=\phi^{-1}((2(i-1)\ell+\ell+\pi_i(j)))=2(i-1)\ell+j=v$.}
%
Indeed, in this case $\{v,w\}$ is an edge in $G_n$
but $\{v,w\}$ is not an edge in $\imu(G_n)$.
%
(Hence, if the number of vertices of this case
is $\Omega(|\{u\in[n]:\mu(u)\neq u\}|)$,
then the difference between $G_n$ and $\imu(G_n)$
is $\Omega(|\{u\in[n]:\mu(u)\neq u\}|)$, and the same holds
with respect to the difference between $\mu(G_n)$ and $G_n$.)
\EDes
Combining all these cases, we get a total contribution
that is proportional to $|\{v\in[n]:\mu(v)\neq v\}|$,
where we might have counted the same contribution
in several different cases.
Since the number of cases is a constant, the theorem follows.
\EPF


\paragraph{Digest: Using large collections of pairwise far apart permutations.}
The construction presented in the proof of Theorem~\ref{construction:thm}
utilizes a collection of $(\ell!)^{\Omega(1)}$ permutations over $[\ell]$
that are pairwise far-apart (i.e., every two permutations differ
on $\Omega(\ell)$ inputs). Such a collection is constructed
in $\tildeO(\ell!)$-time by an iterative exhaustive search,
where the permutations are selected iteratively such that
in each iteration we find a permutation that is far from
permutations that were included in previous iterations.
We mention that in Section~\ref{step3:sec} we shall use
a collection of $\exp(\Omega(\ell))$ such permutations that is
locally computable (i.e., given the index of a permutation
we find its explicit description in polynomial time).
We also mention that, in follow-up work~\cite{GW:perm},
we provided a locally computable collection of $(\ell!)^{\Omega(1)}$
that are pairwise far-apart.


\paragraph{Digest: Combining two robustly self-ordered graphs.}
One ingredient in the proof of Theorem~\ref{construction:thm}
is forming connected components that consist of two
robustly self-ordered graphs that have different vertex degrees
and are connected by a bounded-degree bipartite graph.
Implicit in the proof is the fact that
the resulting graph is robustly self-ordered graph.

\BCM[Combining two $\Omega(1)$-robustly self-ordered graphs]
\label{bd:bipartite:clm}
%
%
%
%
%
%
%
%
%
%
For $i\in\{1,2\}$ and constant \mbox{$\gamma>0$},
let $G_i=(V_i,E_i)$ be an $\gamma$-robustly self-ordered graph,
and consider a graph $G=(V_1\cup V_2,E_1\cup E_2\cup E)$
of maximum degree $d$
such that $E$ contain edges with a single vertex in each $V_i$;
that is, $G$ consists of~$G_1$ and $G_2$
and an arbitrary bipartite graph that connects them.
If the maximum degree in~$G$ of each vertex in $V_1$
is strictly smaller than the minimum degree of each vertex in $V_2$,
then~$G$ is $\gamma/(4d+1)$-robustly self-ordered.
\ECM

\BPFS
For an arbitrary permutation $\mu:V\to V$,
let $T$ denote the set of its non-fixed-points,
and consider the following two cases.

\BDes
\item[{\em Case 1}:] More than $t=\gamma'\cdot|T|$ vertices are
mapped by $\mu$ from~$G_1$ to $G_2$, where $\gamma'=\gamma/(4d+1)$.

In this case,
we get a contribution of at least one unit per each such vertex,
due to the difference in the degrees between $V_1$ and $V_2$.

\item[{\em Case 2}:] at most $t$ vertices are
mapped by $\mu$ from~$G_1$ to $G_2$.

In this case, letting $T_i$ denote the set of non-fixed
vertices in $G_i$ that are mapped by $\mu$ to $G_i$,
we get a contribution of
at least $\sum_{i=1,2}(\gamma\cdot|T_i|-d\cdot t)$ units,
where the negative term is due to possible change in the incidence
with vertices in $T\setminus T_i$.
Hence, the total contribution in this case is at least
$\gamma\cdot(|T|-2t)-2d\cdot t \geq \gamma\cdot|T|-4dt= \gamma'\cdot|T|$.
%
%
%
\EDes
The claim follows.
\EPFS




\paragraph{Regaining regularity and expansion.}
While Theorem~\ref{construction:thm} achieves our main objective,
it useful towards some applications
(see, e.g., the proof of Theorem~\ref{strong-construction:thm})
to obtain this objective with graphs that are both regular and expanding.
This is achieved by applying Theorem~\ref{make:regular+expanding:thm}.
Hence, we have.

\BT[Theorem~\ref{construction:thm}, revised]
\label{construction-add:thm}
For any sufficiently large constant~$d$,
there exists an efficiently constructable family $\{G_n\}_{n\in\N}$
of robustly self-ordered $d$-regular expander graphs.
That is, there exists a polynomial-time algorithm
that on input $1^n$ outputs the $n$-vertex graph $G_n$.
\ET

\subsection{Strong (i.e., local) constructions}
\label{step3:sec}
%
While Theorem~\ref{construction-add:thm} provides
an efficient construction of
robustly self-ordered $d$-regular expander graphs,
we seek a stronger notion of constructability.
Specifically, rather than requiring that the graph be constructed
in time that is polynomial in its size, we require that the neighbors
of each given vertex can be found in time that is polynomial
in the vertex's name
(i.e., time that is polylogarithmic in the size of the graph).
We call such graphs {\sf locally constructable}
(and comment that the term ``strongly explicit''
is often used in the literature).

\BT[Locally constructing robustly self-ordered graphs]
%
\label{strong-construction:thm}
For any sufficiently large constant~$d$, there exists
a locally constructable family $\{G_n=([n],E_n)\}_{n\in\N}$
of robustly self-ordered $d$-regular graphs.
That is, there exists a polynomial-time algorithm
that on input $n$ and $v\in[n]$ outputs the list of neighbours
of vertex $v$ in $G_n$.
Furthermore, the graphs are either expanders
or consist of connected components of logarithmic size.
%
\ET
(Indeed, this establishes Theorem~\ref{main:ithm}.)
We comment that using the result of~\cite{GW:perm},
we can also get connected components of sub-logarithmic size,
as in Theorem~\ref{construction:thm}.
\footnote{Specifically, the result of~\cite{GW:perm} provides
a construction of a collection of $L=\exp(\Omega(\ell\log\ell))$
permutations over $[\ell]$ that are pairwise far-apart
along with a polynomial-time algorithm that, on input $i\in[L]$,
returns a description of the $i^\xth$ permutation
(i.e., the algorithm should run in $\poly(\log L)$-time).
Using this algorithm,
we can afford to set $\ell=\frac{O(\log n)}{\log\log n}$
as in Theorem~\ref{construction:thm}.}
\medskip

\BPF
We employ the idea that underlies the proof of
Theorem~\ref{construction:thm},
while starting with an {\em efficiently constructable family}\/
of robustly self-ordered graphs
(as provided by Theorem~\ref{construction-add:thm})
rather than with the mere existence of a family of such graphs
(equiv., with $\ell$-vertex graphs that can be constructed
in $\poly(\ell!)$-time).
We use a slightly larger setting of $\ell$,
which allows us to use a collection of $\exp(\Omega(\ell))$
pairwise-far-apart permutations (rather than a collection
of $\exp(\Omega(\ell\log\ell))$ such permutations).
Lastly, we apply the same transformation as in the
proof of Theorem~\ref{construction-add:thm}
(so to regain regularity and expansion).
Details follow.

Given a generic $n$, let $\ell=O(\log n)$,
which implies that $\exp(\ell)=\poly(n)$.
By Theorem~\ref{construction-add:thm},
for all sufficiently large $d'$,
we can construct $\ell$-vertex $d'$-regular expander graphs that
are robustly self-ordered (with respect to the robustness parameter $c$)
in $\poly(\ell)$-time.
Again, we shall use two such graphs:
a $d'$-regular graph, denoted $G'_\ell=([\ell],E'_\ell)$,
and a $d''$-regular graph, denoted $G''_\ell=([\ell],E''_\ell)$,
where $d''=d'+1$.

Using $G'_\ell$ and $G''_\ell$, we construct an $n$-vertex
robustly self-ordered graph, denoted $G_n$,
that consists of $n/2\ell$ connected components that
are pairwise far from being isomorphic to one another.
This is done by picking $m=n/2\ell$ permutations,
denoted $\pi_1,\ldots,\pi_m:[\ell]\to[\ell]$,
that are pairwise far-apart, and constructing $2\ell$-vertex graphs
%
%
%
such that the $i^\xth$ such graph consist of
a copy of $G'_\ell$ and a copy of $G''_\ell$ that are connected
by a matching as determined by the permutation $\pi_i$
(as detailed in \eqref{edges-of-cc:eq}).
%

Using the fact that $m<2^\ell$
(rather that $m=\exp(\Theta(\ell\log\ell))$),
we can construct each of these permutations in $\poly(\ell)$-time
by using sequences of disjoint traspositions determined
via a good error correcting code.
Specifically, for $k=\log_2m<\log_2n$,
we use an error correcting code $C:\bitset^k\to\bitset^\ell$
of constant rate (i.e., $\ell=O(k)$) and linear distance
(i.e., the codewords are $\Omega(\ell)$ bits apart from each other),
and let $\pi_i(2j-1)=2j-1+C(i)_j$ and $\pi_i(2j)=2j-C(i)_j$,
where $i\in[m]=[2^k]\equiv\bitset^k$ and $j\in[\ell/2]$.
(That is, the $i^\xth$ permutation switches the pair $(2j-1,2j)\in[\ell]^2$
if and only if the $j^\xth$ bit in the $i^\xth$ codeword is~1,
where $C(i)$ is considered the~$i^\xth$ codeword.)

Like in the proof of Theorem~\ref{construction:thm},
the $i^\xth$ connected component of $G_n$ is isomorphic
to a graph with the vertex set $[2\ell]$
and the edge set
\begin{equation}\label{edges-of-cc:eq}
E'_\ell \;\cup\; \{\{\ell+u,\ell+v\}:\{u,v\}\in E''_\ell\}
         \;\cup\; \{\{v,\ell+\pi_i(v)\}:v\!\in\![\ell]\}.
\end{equation}
The key observation is that, for every $i\in[m]$ and $j\in[\ell]$,
the neighborhood of the $j^\xth$ (resp., $(\ell+j)^\xth$) vertex
in the $i^\xth$ connected component of the $n$-vertex graph $G_n$
is determined by $G'_\ell$ and~$\pi_i(j)$
(resp., by $G''_\ell$ and $\pi_i^{-1}(j)$),
which in turn are constructible in $\poly(\ell)$-time.
Hence, the neighborhood of each vertex in $G_n$
can be found in $\poly(\ell)$-time.
This implies local constructability, since $\ell=O(\log n)$.

The fact that $G_n$ is robustly self-ordered was already
established in the proof of Theorem~\ref{construction:thm},
which is oblivious of the permutations used as long as any
pair of permutations disagrees on $\Omega(\ell)$ points.
Lastly, we may obtain regularity and expansion
by applying Theorem~\ref{make:regular+expanding:thm}.
%
%
%
%
\EPF


%
%
\subsection{Local self-ordering}
\label{local-so:sec}
Recall that, by Definition~\ref{asymmetric:def},
a graph $G=([n],E)$ is called {\sf self-ordered}
if for every graph $G'=(V',E')$ that is isomorphic to~$G$
there exists a unique bijection $\phi:V'\to[n]$ such that $\phi(G')=G$.
One reason for our preferring the term ``self-ordered'' over
the classical term ``asymmetric'' is that we envision being
given such an isomorphic copy $G'=(V',E')$ and asked to find its
unique isomorphism to~$G$, which may be viewed as ordering
the vertices of $G'$ according to (their index in) $G$.
The task of finding this unique isomorphism
will be called {\em self-ordering $G'$ according to~$G$}\/
or {\em self-ordering $G'$}\/ (when~$G$ is clear from the context).

Evidently, the task of self-ordering a given graph $G'$
according to a (self-ordered) graph~$G$ that can
be efficiently constructed reduces to testing isomorphism.
When the graphs have bounded-degree the latter task
can be performed in polynomial-time~\cite{L}.
%
These are general facts that do apply also to
the robustly self-ordered graph $G_n$ constructed
in the proof of Theorem~\ref{strong-construction:thm}.
However, in light of the fact that the graph $G_n$
is {\em locally}\/ constructable, we can hope for more.
Specifically, it is natural to ask if we can perform
self-ordering of a graph $G'$ that is isomorphic to $G_n$
in a {\em local}\/ manner; that is, given a vertex in $G'$
(and oracle access to the incidence function of $G'$),
can we find the corresponding vertex in $G_n$ in $\poly(\log n)$-time?
Let us define this notion formally.
%
%
%
%
%


%
\BD[Locally self-ordering a self-ordered graph]
\label{local-so:def}
We say that a self-ordered graph $G=([n],E)$
is {\sf locally self-ordered}
if there exists a polynomial-time algorithm that,
given a vertex $v$ in any graph $G'=(V',E')$ that is isomorphic to~$G$
and oracle access to the incidence function of $G'$,
finds $\phi(v)\in[n]$ for the unique bijection  $\phi:V'\to[n]$
such that $\phi(G')=G$
{\rm(i.e., the unique isomorphism of $G'$ to~$G$)}.
\ED
Indeed, the isomorphism $\phi$ orders the vertices of $G'$
in accordance with the original (or target) graph~$G$.
We stress that the foregoing algorithm works in time that is polynomial
in the description of a vertex (i.e., $\poly(\log n))$-time),
which is polylogarithmic in the size of the graph (i.e., $n$).
%
We show that such algorithms exist for the graphs constructed
in the proof of Theorem~\ref{strong-construction:thm}.

%
%
\BT[Locally self-ordering the graphs
of Theorem~\ref{strong-construction:thm}]
\label{local-so4strong-construction:thm}
For any sufficiently large constant~$d$, there exists
a locally constructable family $\{G_n=([n],E_n)\}_{n\in\N}$
of robustly self-ordered $d$-regular graphs
that are locally self-ordered.
%
%
%
Furthermore, the graphs are either expanders
or consist of connected components of logarithmic size.
\ET
As in Theorem~\ref{strong-construction:thm},
we can obtain connected components of sub-logarithmic size
by using~\cite{GW:perm}.
\medskip

\BPF
We first consider the version that yields $n$-vertex graphs
that consist of connected components of logarithmic size.
The basic idea is that it we can afford reconstructing
the connected component in which the input vertex reside,
and this allows us both to determine the index of the vertex
in this connected component as well as the index of the component
in the graph. Specifically,
on input a vertex $v$ in a graph $G'$ that is isomorphic to $G_n$,
we proceed as follows.

\BE
\item
Using queries to the incidence function of $G'$,
we explore and retrieve the entire $2\ell$-vertex
connected component in which $v$ resides, where $\ell=\log_2n$.

Recall that this connected component consists of (copies of)
two $\ell$-vertex regular graphs, denoted $G'_\ell$ and $G''_\ell$,
that are connected by a matching.
Furthermore, these graphs have different degrees
and are each (robustly) self-ordered.
\item
Relying on the different degrees,
we identify the foregoing partition of this $2\ell$-vertex component
into two $\ell$-vertex (self-ordered) graphs,
denoted $A_v$ and $B_v$, where $A_v$ (resp., $B_v$)
is isomorphic to $G'_\ell$ (resp., $G''_\ell$).
\item
Relying on the self-ordering of $G'_\ell$ (resp., $G''_\ell$),
we order the vertices of $A_v$ (resp., $B_v$).
This is done by constructing $G'_\ell$ (resp., $G''_\ell$),
and using an isomorphism tester.
The order of the vertices in $A_v$ and $B_v$ also determines
the permutation that defines the matching between the two graphs.
\item
Relying on the correspondence between the permutations used
in the construction and codewords of a good error-correcting code,
we decode the relevant codeword (i.e., this is decoding without error).
This yields the index of the permutation in the collection,
which equals the index of the connected component.
\EE
Note that this refers to the basic construction that was presented
in the proof of Theorem~\ref{strong-construction:thm},
before it was transformed to an expander and made regular.
Recall that both transformations are performed by augmenting
the graph with auxiliary edges that are assigned a different color
than the original edges, and that edges with different colors
are later replaced by copies of different (constant-size) gadgets.
These transformations do not hinder the local self-ordering procedure
described above, since it may identify the original graph
(and ignore the gadgets that replace other edges).
The claim follows.
\EPF
%


\paragraph{Local reversed self-ordering.}
While {\em local self-ordering}\/ a (self-ordered) graph
seems {\em the natural local version}\/ of self-ordering the graph,
an alternative notion called {\em local reversed self-ordering}\/
will be defined and studied next
(and used in Section~5).
Both notions refer to a self-ordered graph, denoted $G=([n],E)$,
and to an isomorphic copy of it, denoted $G'=(V',E')$;
that is, $G=\phi(G')$ for a (unique) bijection  $\phi:V'\to[n]$.
While local self-ordering is the task of
finding the index of a given vertex of $G'$ according to~$G$
(i.e., given $v\in V'$, find $\phi(v)\in[n]$),
local reversed self-ordering is the task of
finding the vertex of $G'$ that has a given index in~$G$
(i.e., given $i\in[n]$, find $\phi^{-1}(i)\in V'$).
%
%
%
%
%
%
In both cases, the graph~$G$ is locally constructible
and we are given oracle access to the incidence function of $G'$.
In addition, in the reversed task, we assume that the algorithm
is given an arbitrary vertex in $G'$,
since otherwise there is no hope to hit any element of $V'$.
\footnote{Needless to say, this is not needed in case $V'=[n]$,
which is the case that is used in Section~5.}

%
\BD[Locally reversed self-ordering]
\label{local-reversed-so:def}
We say that a self-ordered graph $G=([n],E)$
is {\sf locally reversed self-ordered}
if there exists a polynomial-time algorithm that,
given $i\in[n]$
and oracle access to the incidence function of
a graph $G'=(V',E')$ that is isomorphic to~$G$
and an arbitrary vertex $s\in V'$,
finds $\phi^{-1}(i)\in V'$ for the unique bijection  $\phi:V'\to[n]$
such that $\phi(G')=G$
{\rm(i.e., the unique isomorphism of $G'$ to~$G$)}.
\ED
We stress that the foregoing algorithm works in time that is polynomial
in the description of a vertex (i.e., $\poly(\log n))$-time),
which is polylogarithmic in the size of the graph (i.e., $n$).
We show that such algorithms exist for variants of the graphs
constructed in the proof of Theorem~\ref{strong-construction:thm}.
In fact, we show a more general result that refers
to any graph that is locally self-ordered and for which
short paths can be locally found between any given pair of vertices.

%
%
\BT[Sufficient conditions for
locally reversed self-ordering of graphs]
\label{local-reversed-so:thm}
Suppose that $\{G_n=([n],E_n)\}_{n\in\N}$ is a family
of bounded degree graphs that is locally self-ordered.
Further suppose that given $v,u\in[n]$,
one can find in polynomial-time a path from $u$
to $v$ in $G_n$.
Then, $\{G_n=([n],E_n)\}_{n\in\N}$
is locally reversed self-ordered.
\ET
We mention that any family of robustly self-ordered graphs
that is locally self-ordered can be transformed
into one that also supports locally finding short paths.
This is done by superimposing the graphs of this family
with graphs that supports locally finding short paths,
while using different colors for the edges of the two graphs
and later replacing these colored edges by gadgets
(as done in Section~\ref{edge-colored:trans:sec}).
We also mention that applying degree reduction to the hyper-cube
(i.e., replacing the original vertices with simple cycles)
yields a graph that supports locally finding short paths.
\footnote{For any $\ell\in\N$, the resulting graph consists of the
vertex-set $\{\ang{x,i}\!:\!x\!\in\!\bitset^\ell\,\&\,i\!\!\in\![\ell]\}$
and edges that connect $\ang{x,i}$ to $\ang{x\xor0^{i-1}10^{\ell-i},i}$
and to $\ang{x,i+1}$, where $\ell+1$ stands for~1.
For simplicity of exposition, we also add self-loops on all vertices.
%
Then, given $\ang{x,i}$ and $\ang{y,j}$, we can combine the $2\ell$-path
that goes from $\ang{x,i}$ to $\ang{y,i}$ with the $|j-i|$-path
that goes from $\ang{y,i}$ to $\ang{y,j}$,
where the odd steps on the first path move from $\ang{z,k}$
to $\ang{z\xor0^{i-1}10^{\ell-i},k}$ (or stay in place)
and the even steps (on this path) move from $\ang{z,k}$ to $\ang{z,k+1}$.}
%
\medskip

\BPF
On input $i\in[n]$ and $s\in V'$,
and oracle access to the incidence function
of a graph $G'=(V',E')$ that is isomorphic to $G_n$,
we proceeds as follows.
\BE
\item
Using the local self-ordering algorithm,
we find $i_0=\phi(s)$, where $\phi:V'\to[n]$
is the unique bijection satisfying $\phi(G')=G$.

\item
Using the path-finding algorithm for~$G$,
we find a $\poly(\log n)$-long path from $i_0$ to $i$ in~$G$.

Let $\ell$ denote the length of the path,
and denote its intermediate vertices by $i_1,\ldots,i_{\ell-1}$;
that is, the full path is $i_0,i_1,\ldots,i_{\ell-1},i_\ell=i$.
\item
Using the local self-ordering algorithm (and our oracle access to $G'$),
we iteratively find the corresponding vertices in $G'$,
denoted $v_1,....,v_\ell$, where $v_j=\phi^{-1}(i_j)$.

For $j=1,\ldots,\ell$, we find $v_j\eqdef\phi^{-1}(i_j)$ as follows.
First, using queries to the incidence function of $G'$,
we find all neighbors (in $G'$) of $v_{j-1}$,
where $v_0\eqdef s$ (and, indeed, $v_0=\phi^{-1}(i_0)$).
Next, using the local self-ordering algorithm,
we find the indices of all these vertices in~$G$; that is,
for every vertex $w$ that neighbors $v_{j-1}$, we find $\phi(w)$.
Last, we set $v_j$ to be the neighbor that has index $i_j$ in~$G$;
that is, $v_j$ satisfies $\phi(v_j)=i_j$.
\EE
Hence, $v_\ell$ is the desired vertex;
that is, $v_\ell$ satisfies $\phi(v_\ell)=i_\ell=i$.

Assuming that
the local self-ordering algorithm has query complexity $q(n)$,
that the paths found in~$G$ have length at most $\ell(n)$,
and that $d$ is the degree bound,
the query complexity of our reversed self-ordering algorithm
is $(1+\ell(n)\cdot d)\cdot(q(n)+1)$,
where we count both our direct queries to the incidence function of~$G$
and the queries performed by the local self-ordering algorithm.
Similar considerations apply to the time complexity.
\EPF


\BCR[A version of Theorem~\ref{local-so4strong-construction:thm}
supporting local reversed self-ordering]
\label{local-reversed-so4rso:cor}
For any sufficiently large constant~$d$, there exists
a locally constructable family $\{G_n=([n],E_n)\}_{n\in\N}$
of robustly self-ordered graphs of maximum degree $d$ that are
both locally self-ordered and locally reversed self-ordered.
\ECR
The corollary follows by combining
Theorem~\ref{local-so4strong-construction:thm}
with Theorem~\ref{local-reversed-so:thm},
while using the augmentation outlined following
the statement of Theorem~\ref{local-reversed-so:thm}.
We mention that Corollary~\ref{local-reversed-so4rso:cor}
will be used in Section~5.


%
%
%
%
%
%
%
%
%
%
%
%
%
%
%
%
%
%



%











































































































































































































































































































































































































































































































































































































































































































































































































































































































































































































































































































































































































































































































































































































































































































































































































































































































































































































































































































































































































































































































































































































































































































































































































































































































































































































































































































































































































































































































































































































































































































































































































































































































































































































































































































































































































































































































































































































































































































































































































































































































































































































































































































































































































































































































































































































































































































